% Derived from the owned full Markdown narrative; source-stage only.


\section{Assistance and Learning}

Imagine watching a hand play a short tune. Every key lands at the right moment. It looks like skill. Now imagine that an outside computer chose the notes and controlled each finger. The tune is still real as an action, but we have not yet shown that the nervous system learned to play it.\par

This is the question behind Brain2Body: could useful help leave a nervous system more capable after the help ends? That is a bigger question than making a body move once. It asks about what was learned, what support is still needed and whether the improvement survives tomorrow, a different starting position or an unfamiliar problem.\par

Brain2Body means brain to body; the 2 stands for ``to.'' Information also travels the other way. A movement changes what the body senses. That new information helps guide the next movement. The programme proposes studying these relationships with real measurements, computer models of bodies and carefully limited experiments. It does not claim that a complete skill can already be downloaded into someone.\par

The long-term ambition is deliberately bold: help a trainee acquire difficult physical skills much faster. Sewing, playing violin, solving a puzzle with the hands or coordinating a dance would be demanding tests. They require different body parts, information and kinds of feedback. These examples set the ambition. They are not a list of things the current system can do.\par

NeuroTranslation proposes mapping nervous-system signals between clearly described situations, such as a particular recording and a particular movement task. Brain2Body takes the next question: what happens when that prediction has to work through a physical body and help a trainee? A useful prediction can start the experiment, but the body's response and the trainee's later performance determine what it has achieved. We need to follow the information into the movement, then follow the movement beyond the moment of help. Each step gives the next test a reason to exist.\par

\section{A Single-Key Test}

Suppose the task is to press the leftmost of four keys, then release it
before the next beat. Success needs several checks: the right key moves, it
moves at the right time, it is released, and a neighbouring key is not pressed
accidentally. State what happens if the finger misses completely. Each rule
names a different way the movement can succeed or fail.\par

That small example contains several problems. The trainee may choose the wrong key. The finger may move toward the right key but fall short. The finger may press correctly and hold too long. A second finger may strike a neighbour. A single label called ``success'' can hide all of these differences unless the rules are stated.\par

A camera can help measure where a finger went. A contact sensor can record that a key was touched. A force sensor can measure how hard it was pressed. A clock can record timing. An electrode can record nearby electrical activity. These are different observations. Looking at a brain signal does not automatically tell us the force at a fingertip; seeing a key press does not automatically tell us which neural mechanism produced it.\par

Before building the model, decide which body and task it is meant to describe, which recording it will receive and which independent measurement will check its answer. Set the allowed error as part of that question. A score can then tell us whether the answer was useful under those conditions. Adding a more elaborate display before defining success leaves the same unanswered question in a larger picture.\par

\section{Movement and Feedback}

Movement is usually a conversation with the world. You reach for a cup. Your hand's position changes. You see the change and feel contact. If the cup slips, you adjust. The action affects the next information you receive, and that information affects the next action.\par

This is a feedback loop. Feedback means information about what happened, returned to help guide what happens next. Feedback control means adjusting action in response to the difference between the target and the observed result. The established work of Todorov and Jordan helps explain movement control under uncertainty. Our proposal uses that idea as a reason to test the whole loop rather than celebrate only a command. \href{https://www.nature.com/articles/nn963}{Todorov and Jordan, 2002}.\par

Timing matters. Information arriving after a movement has finished may still help the next attempt, but it did not correct the movement already completed. A recording played back later is not live feedback to the person who produced it. Feedback generated by a computer body is not measured feedback from a living body. Those differences should remain visible beside the result.\par

``Perception is subjective. Movement is not'' is the programme's starting contrast. A person experiences a sensation. Other people can independently measure the body's position, trajectory and forces. That does not mean measurements are perfect, or that we can measure the private feeling directly. It means movement offers an outcome we can check without pretending to read someone's entire inner life.\par

If a model predicts that a finger should move five millimetres, a suitable measurement can compare its prediction with the actual movement. The error becomes information. Did the model overshoot? Did it start late? Did it choose a different finger? These questions help locate the failure and guide the next experiment. They are more useful than an unexplained claim that the model is ``brain-like.''\par

\section{Accuracy and Repeatability}

Imagine a ruler whose markings are wrong. Each time you measure the same object, you get almost the same answer. The measurement is repeatable, but it may be inaccurate. A computer model can have the same problem: a confident, stable answer can still disagree with the independent reference.\par

An accurate measurement agrees with an independent reference under declared conditions. A precise measurement repeats consistently under stated conditions. Ground truth is that independently acquired reference, together with its collection method, calibration and limits. Calling it ground truth does not make it perfect. \href{https://www.nist.gov/pml/nist-technical-note-1297/nist-tn-1297-appendix-d1-terminology}{NIST measurement terminology}.\par

Calibration means checking how an instrument's readings relate to the quantity we want to measure. A sensor that reports values between zero and one may describe relative finger movement. Without the necessary calibration, it does not automatically tell us millimetres or a joint angle. A model can correctly predict those relative values while still leaving the physical body question unanswered.\par

More detail is not a shortcut around this problem. A simulation can update thousands of times a second. A display can show many decimal places. Neither proves that the underlying values match biology. Useful claims show the predicted and measured outcomes, the error, repeated results and the conditions where the comparison holds.\par

\section{Structure and Function}

A road map shows connections. It does not show today's traffic, an unexpected road closure or why a particular driver chose a turn. The map is valuable because it constrains what routes are possible. It is not a complete explanation of every journey.\par

A map of neural connections has a similar role. It can show which cells connect and where. It does not by itself tell us how strongly one affects another in a living subject at a particular moment. Activity depends on timing, chemicals, bodily state, past learning and the task being performed. The programme needs both structural constraints and suitable observations of function.\par

The Human Genome Project is a useful comparison. Reading a reference sequence did not finish the investigation of what its parts do or how they are regulated. Further research investigated functional elements. That does not mean the sequencing project failed. It means obtaining a map and explaining function are different achievements. \href{https://www.genome.gov/11009066/2003-release-scientists-venture-deeper-into-the-human-genome}{NHGRI, Beyond Genes}.\par

MICrONS is an important example of bringing different kinds of evidence together. Its researchers related recorded visual responses to detailed connections in a region of one mouse's brain. That helps ask how structure and observed response relate. It does not supply a complete human brain, a complete body or every nerve pathway needed to teach a new skill. \href{https://www.nature.com/articles/s41586-025-08790-w}{MICrONS Consortium, 2025}.\par

The map's coverage determines what we can ask of it. A map of one brain region does not cover the whole brain, and a brain map does not cover the whole nervous system. Relating an electrode to a particular cell in an atlas needs evidence that both refer to that cell in the recorded subject. The next question is therefore about the recording itself: what did the instrument observe, and what must we infer from it? A detailed map cannot answer on the instrument's behalf.\par

\section{What Brain Measurements Record}

Functional magnetic resonance imaging, usually shortened to fMRI, measures changes in a blood-related signal. It takes measurements over time. It does not directly photograph thoughts. To interpret it, researchers need an account of how neural activity, blood flow and the instrument's readings relate.\par

Logothetis and colleagues studied these relationships. In their experiments, the blood-related signal was more closely related to local field potentials than to spiking output. Local field potentials are an aggregate electrical measurement near a recording site; spikes are brief electrical events associated with neurons. Those are distinct quantities, and the relationship varied with the recording conditions. \href{https://pubmed.ncbi.nlm.nih.gov/11449264/}{Logothetis et al., 2001}.\par

Tang and colleagues showed that a system could reconstruct aspects of meaning from particular brain recordings under constrained conditions. The result depended substantially on the person and cooperation. It does not establish a universal decoder that can recover every person's private experience word for word. \href{https://www.nature.com/articles/s41593-023-01304-9}{Tang et al., 2023}.\par

The language-translation analogy is helpful here. Words can mean different things in different contexts; a translation may preserve some meaning while leaving ambiguity. Neural translation must likewise state what information is preserved and where the match is uncertain. The analogy does not prove that useful common representations are impossible. It prevents us from assuming a perfect one-to-one correspondence before testing it.\par

\section{Measurement Effects and the Quantum Analogy}

Sometimes observing something changes the conditions under which it behaves. A participant in a scanner may have to remain still and follow an unusual task. An implanted interface may affect tissue. Those practical interactions need to be measured and explained.\par

Quantum measurement can provide an analogy about observation and interaction. But an analogy is not a demonstrated brain mechanism. Individual quantum systems and large-scale neural measurements require different explanations. Quantum backaction means that a measurement changes the state of the quantum system being measured. Quarks are subatomic particles found inside protons and neutrons. We do not claim that these effects or measurements at that scale explain the limits of fMRI or the difficulty of translating subjective experience. \href{https://www.nobelprize.org/uploads/2018/06/advanced-physicsprize2012-1.pdf}{Royal Swedish Academy, quantum measurement background}.\par

The limits already have concrete explanations to investigate: an indirect signal, limited sampling, task constraints, differences between individuals and missing measurements. Each points toward a particular check and a reason to state uncertainty. Calling the problem ``quantum'' would supply none of those missing measurements or make a prediction more accurate.\par

\section{Existing Brain--Body Interfaces}

There are published examples of brain signals helping produce real movement. Bouton and colleagues recorded from motor brain tissue and used the information to control electrical activation of forearm muscles. A participant with quadriplegia, paralysis affecting all four limbs, performed useful hand and wrist actions. That is a real achievement with a defined participant and interface. It is not evidence that an unfamiliar skill was uploaded. \href{https://www.nature.com/articles/nature17435}{Bouton et al., 2016}.\par

Lorach and colleagues connected recorded brain signals to spinal stimulation in a participant with chronic spinal injury. They evaluated standing and walking in the person, and reported rehabilitation-associated improvement with the interface off. A body actually moving under a measured intervention provides a stronger functional test than a computer animation alone. The result remains tied to that person, intervention and rehabilitation history. \href{https://www.nature.com/articles/s41586-023-06094-5}{Lorach et al., 2023}.\par

Flesher and colleagues supplied touch-related sensations by stimulating sensory brain tissue during robotic arm control. The feedback helped the participant control the arm. This supports asking whether useful information returning toward the brain improves action. It does not show that a whole expert skill transferred from another person. \href{https://pubmed.ncbi.nlm.nih.gov/34016775/}{Flesher et al., 2021}.\par

The movement bridge and the return of touch give this programme two directions to investigate. They were tested in separate participants and tasks. A new combined system must therefore show how its own measurements and timing fit together, then report its outcomes and failures. Reading the papers together gives us a reason to build that experiment; it cannot supply evidence for links that the combined experiment has not yet tested.\par

\section{Digital-Fly Models}

The digital fly makes this problem harder to dodge. A fly on screen walks;
we want to know what made it walk. Was the decisive ability supplied by its
connection map, the equations chosen for its neurons, a trained body controller
or information added from outside? These can all be useful components. Keeping
their jobs visible tells us which biological claim the result can actually
test.\par

FlyWire mapped 139,255 neurons in one adult female fly brain. A neuron is a nerve cell; a connectome is a map of neural connections. The resource is powerful, but it also has reconstruction errors and missing information. Its chemical connections do not directly measure every physiological effect. It is not the entire nervous system, body and environment. \href{https://www.nature.com/articles/s41586-024-07558-y}{Dorkenwald et al., 2024}.\par

Pospisil and colleagues propose a way to use connection maps together with interventions to estimate causal effects. An intervention deliberately changes something to test what it affects. Their estimator tests use simulated conditions where the underlying answer is known by construction. Those tests do not mean the full causal dynamics of a living fly have already been measured. \href{https://www.nature.com/articles/s41586-024-07982-0}{Pospisil et al., 2024}.\par

Shiu and colleagues show how a useful approximation can earn evidence. They made particular predictions from a connectome-based model, then compared them with experiments involving stimulation and behaviour. Many predictions agreed. Others did not. Experimentally observed effects associated with cells called Phantom and Usnea were missed, and water-related predictions for G2N-1/Usnea disagreed with silencing tests, which reduce the activity of the targeted cells. A cell's name does not explain away a mismatch. \href{https://pmc.ncbi.nlm.nih.gov/articles/PMC11446845/}{Shiu et al., 2024}.\par

The researchers reported agreement for 91\% of 164 tested predictions, or 84\% after excluding one mostly negative experimental screen. These percentages describe those predictions and that comparison. They do not say the model is 91\% of a complete fly. Missing receptor dynamics, electrical connections and chemical regulation matter to the explanation, including why absolute simulated firing rates are not expected to be fully accurate.\par

A model earns support prediction by prediction. It can suggest a useful experiment, agree with that experiment and still fail in another condition. Reporting the limits beside the success shows both what the result supports and which further measurement could test the explanation. That is how an approximation can contribute without being mistaken for a complete animal.\par

\section{Foot Contact and Simulation Errors}

A detailed body model must solve numerical problems as well as anatomical ones. NeuroMechFly v2 reports that simulated leg tips can sit slightly below the ground because of numerical errors or physics instability. The authors used a 0.05 millimetre margin to avoid detecting a false overstretch. That is a specific contact artifact documented in the paper. \href{https://gizemozd.github.io/assets/pdf/2024_neuromechflyv2.pdf}{Wang-Chen et al., 2024}.\par

The reported failure is a small below-ground position and the modelling response to it. It supports a precise claim about contact in that simulator, not a claim that a FlyWire brain acquired an ability no living fly could ever perform. There is no demonstrated gross tunnelling, enormous jump or perpetual motion in this example. Naming the actual failure lets us investigate the correction instead of expanding it into an unsupported story.\par

The same paper models adhesion with engineered forces and defined stepping phases while the relevant biological liftoff mechanisms remained uncertain. The computer can perform a movement using that substitute. Whether the substitute reproduces the animal's mechanism is another question. An unknown part of biology does not become known because software needs a value there.\par

FlyBody provides another example. Its artificial controllers were trained to imitate recorded walking and flight. That can test how well a model matches those measured movements. It does not recover a FlyWire brain. The paper's time-averaged treatment of aerodynamics and its missing detailed muscle and sensory physiology remain limitations. Similar movement can be produced by different internal mechanisms. \href{https://pmc.ncbi.nlm.nih.gov/articles/PMC12310536/}{Vaxenburg et al., 2025}.\par

\section{Encoders, Controllers and Added Information}

Eon's March 2026 account describes an integrated digital fly using neural and body models. The developers disclose hand-chosen brain--body mappings, substantial dependence on existing body controllers, largely missing internal state and learning mechanisms, and neural dynamics not yet validated against important biological signatures. That is evidence about the disclosed engineered system at that date. It is not proof of a complete biological upload. We have not independently reproduced its dynamics. \href{https://eon.systems/updates/embodied-brain-emulation}{Eon Systems, 10 March 2026}.\par

Games make the attribution problem especially clear. In a pinned fly-chess project, inspected code supplies a symbolic description of the board, selected mathematical dynamics, a trained output reader and outside legal-move search. Those pieces can all contribute to the software's behaviour. Their contribution cannot be credited to fly wiring alone. The project's playing strength was not independently reproduced here. \href{https://github.com/mncrftfrcnm/making-fly-play-chess/blob/b97dfefbf9b8850f3cdf6b412c2ed840be94e260/scripts/fly_chess_trainer.py}{Pinned fly-chess trainer}.\par

DoomFly uses a male nervous-system map, distinct from the female FlyWire brain. Its pinned code reads enemy direction and distance from game-engine geometry and supplies that information to model neurons. The authors report that the intended visual pathway failed in their tests. This is an engineered shortcut to information, not a measurement showing fly vision or a living fly learning Doom. Historical performance remains unreproduced here, with a reported checkpoint absent. \href{https://github.com/shreyash-sharma/doomFly/blob/5f6c2881be5f5f40b0592a75e6ee231b50333b40/src/doom/fearless/danger.py}{Pinned DoomFly input}.\par

A combined software system may be useful because its parts supply different capabilities. An outside planner may choose legal moves from a known board; a body controller may supply locomotion. Disclose those contributions so the reader can decide what the neural model itself explains. Combining them does not make their engineered abilities evidence about the living fly.\par

\section{Reward and Task Performance}

Once we know where the commands came from, we still have to ask what success
meant. A score pays the optimizer for whatever the score measures. If finishing
the task and collecting the reward come apart, the computer can get better at
the latter while leaving the former undone.\par

An optimizer is the procedure that changes a model to improve a chosen objective. If the objective rewards the wrong thing, improvement can be misleading. CoastRunners is a documented example: an agent repeatedly circled reward targets instead of finishing the race. Its score improved while the intended outcome remained unmet. That is a problem with the reward, not proof of impossible physics. \href{https://deepmind.google/blog/specification-gaming-the-flip-side-of-ai-ingenuity/}{Google DeepMind, specification gaming}.\par

Historical first-hand accounts describe a different class of failure. Simulated creatures found movement that exploited numerical errors. Some virtual soft robots penetrated the ground between updates and received corrective forces that supplied artificial energy. Their movement was impressive in the simulator but did not reveal a physically realizable mechanism. Fixing the error changed what they could do. These examples concern particular simulators, not a demonstrated FlyWire defect. \href{https://arxiv.org/html/1803.03453v4}{Lehman et al., first-hand accounts}.\par

For Brain2Body, the response is to measure what the task truly requires. Did the right key move, at the right time, without an extra press? Were contact forces and joint limits physically supported? Did the model receive information the real organism would lack? A high reward cannot answer all of those questions by itself.\par

We also need to distinguish an actual violation from an unfamiliar behaviour. A surprising movement is not automatically impossible. To make that claim, we would need an exact model version, the trajectory and interface, independently supported biological or physical bounds, and a reproducible violation. We currently have no verified evidence here that a FlyWire-derived brain performs something no living fly could ever perform.\par

\section{Protein Structure and Biological Function}

AlphaFold predicts aspects of molecular structure. A convincing shape is not the same thing as a protein shown to bind or function as intended in a living setting. The reference measurement and the tested endpoint matter. The original AlphaFold3 paper describes plausible invented structure in unstructured regions and, separately, predicted chains that overlap. Those are documented structural problems, not a reason to dismiss every output. \href{https://www.nature.com/articles/s41586-024-07487-w}{Abramson et al., 2024}.\par

One result makes the scope issue concrete. The paper reports 62.9\% cluster-averaged success for protein--antibody interfaces at DockQ > 0.23. DockQ is a score comparing a predicted interface with an experimentally determined one; 0.23 is the declared cutoff here. The complement, 37.1\%, failed that threshold. The benchmark contained 65 interface clusters drawn from 71 complexes and used confidence-based selection from 1,000 seeds with five samples per seed for each target. It does not say that 37.1\% of every protein prediction fails to function. Structure and function are different questions.\par

Crystallization can make measurement difficult, but it is not the only method for determining structures. Cryo-electron microscopy uses electrons to study samples; nuclear magnetic resonance uses the response of atomic nuclei to magnetic conditions. Those methods have their own limitations. We should not attribute every unknown structure to difficulty making crystals. \href{https://pdb101.rcsb.org/learn/guide-to-understanding-pdb-data/methods-for-determining-structure}{RCSB structure methods}.\par

Generated examples can help training. AlphaFold2 used predicted structures for some unlabelled sequences to improve learning. But a generated label is not an independent biological measurement. If two models agree because one learned from the other, their agreement cannot supply the missing real-world test. Synthetic data, meaning generated rather than directly measured data, can support a hypothesis while leaving its empirical validation open. \href{https://www.nature.com/articles/s41586-021-03819-2}{Jumper et al., 2021}.\par

\section{Independent Prediction Tests}

Biological virtualization is a demanding claim: a digital model has been shown through independent measurements to reproduce specified biological behaviour and causal dynamics in a stated domain, within stated error limits. A lifelike animation, a copied connection map or a high score does not establish it.\par

Hallucination, in this programme's usage, means invented or ungrounded content presented as reality. If software invents a physiological mechanism and its output is presented as established biological fidelity, that is the problem. An openly disclosed approximation is different. It can be useful, testable and valid for particular outcomes.\par

Suppose a model gets one movement right. That success is worth recording. It does not alone establish that its internal explanation is faithful, or that it will be reliable in a new condition. The dice comparison points to uncertainty about unverified reliability. It does not mean every prediction is random or equally likely. Learned patterns can help make a prediction; its performance still needs to be evaluated.\par

Independent measurements and held-out tests show where a prediction works. Measured errors, failure cases and checked confidence show how much trust that result warrants. The same model may be adequate for a stated task and inadequate for a stronger biological claim; defining that boundary lets a reader use the result without inheriting claims it has not earned.\par

\section{Recording Access, Context and Stable Signals}

Can we obtain a signal that helps answer the question? ``Access it'' begins
there: where is the recording made, what does it measure, and what changes when
the tissue or subject is observed? More electrodes can enlarge the data file
without supplying the missing information. The useful recording must help
distinguish explanations.\par

An available signal then needs a relationship we can check. ``Ground it'' asks
how it relates to independently measured body behaviour. A label saying ``index
finger'' could mean an attempted action, a sensor, a simulated goal or an
anatomical region. The actual meaning, clock and checking measurement determine
which relationship is being tested; the shared words do not make the quantities
interchangeable.\par

``Keep it'' asks whether that recording and mapping stay useful over time.
Tissue may respond to an interface, the recorded population may change, or
calibration and task conditions may shift. A relationship that works on one
date could fail on another. Repeated measurements and appropriate tests must
show what remained usable.\par

Keeping a signal usable is not the same as keeping a learned skill. An electrode may remain stable while a trainee gains no independent capability. A trainee may improve while a recording becomes unreliable. Brain2Body treats these as separate outcomes because each can succeed or fail without the other.\par

\section{Connecting Measurements Across the Body}

Electrical recordings can come from different places. Electroencephalography, or EEG, uses scalp electrodes. Electrospinography, or ESG, records over the spine. Electroneurography, or ENG, records nerve electrical activity. Electromyography, or EMG, records muscle electrical activity. EMG is not direct recording from a peripheral nerve, the nerves outside the brain and spinal cord.\par

The same screen can show all four types. That does not make them simultaneous measurements from the same voluntary action. Different subjects, dates, tasks, instruments and clocks must remain different. Otherwise the display can suggest a chain that the evidence never recorded.\par

The programme's full objective includes brain, spinal, peripheral and body relationships. Open resources from HHMI/Janelia, Allen, Harvard and other groups can support different parts of that work. Structural maps, activity recordings and movement observations have different roles. A wide search does not mean every necessary link has been acquired or paired.\par

When a link is missing, a model may predict it. That prediction could be worth testing. It stays a prediction until an appropriate independent observation checks it. Putting a predicted spinal signal between measured cortex and measured movement cannot silently turn the middle of the chain into measured physiology.\par

\section{Simulation-Based Error Analysis}

MyoSuite provides research simulations where muscles produce movement and objects can make contact. That makes it useful for asking how commands lead to physical outcomes. The tool's existence does not make every simulated hand an accurate copy of the recorded person's hand. \href{https://proceedings.mlr.press/v168/caggiano22a.html}{Caggiano et al., 2022}.\par

Think of four steps. A translator predicts a limited goal from a recording. A controller chooses muscle commands. The body model determines the resulting movement and contact. The task rules score the outcome. Keeping the steps visible lets us see where a failure began.\par

If the translator chooses the middle finger when the label was index finger, a competent controller may press beautifully with the middle finger. The controller executed its command. The complete task still failed. If the predicted finger is correct but no contact occurs, the controller or body model may be responsible. Those explanations lead to different next tests.\par

A direct animation can hide the distinction. Setting a joint to the desired angle is a way to display a pose, but it bypasses the problem of choosing muscle commands that produce it. If the research claim concerns muscle-driven control, the movement must come from the tested controls, and its physical state must be traceable.\par

NeuroTranslator is the programme's name for software joining particular recordings, predictions, replay and computer-body checks. NeuroTranslation Network is a proposed design for joining more components while preserving the origin and meaning of their values. The names describe software and a plan. They do not mean every nerve pathway has been measured.\par

\section{The Six-Session Macaque Study}

The current development study uses six recordings from one macaque: four training dates and two development-validation dates. Training material is used to adjust a model. Development validation helps examine or select a method during development. It is different from untouched material reserved for independent confirmation.\par

Ninety-six contacts in the source's HandM1 motor-cortex region supply two kinds of measurement. Threshold-crossing-rate features count detected signal crossings under the recording method. Spiking-band-power features describe signal power in an archived frequency band. Their short labels are TCR and SBP. Together they give 192 features from 96 contacts, not 192 separately located electrodes.\par

The targets are two source-normalized finger-sensor groups. One represents index-finger movement; the other combines middle, ring and small fingers. ``Source-normalized'' means the source values have been scaled under the declared processing. These values are not a complete set of finger joint angles, forces or body positions in physical units.\par

This is a bounded hand-mapping study. It does not establish a whole-body neural translator, transfer across species or a way to teach a living nervous system. Those larger aims remain. They have not become achievements because the smaller study has finished its training runs. \href{research-evidence.json}{Programme computational evidence summary}.\par

\section{Current Study: Neither Model Passed}

Eighteen runs saved a total of 720 epochs. An epoch is a pass through the training material. Finishing that work is an execution result. It does not tell us whether the model met the performance rules.\par

Both candidates use a gated recurrent unit, usually shortened to GRU. This sequence model can retain information from earlier inputs. The comparison includes a ridge model, a simpler predictor whose training discourages overly large fitted coefficients, the weights assigned to inputs. The question is whether the more elaborate method contributes a benefit under the same rules; its name cannot answer that comparison.\par

Both GRU model gates FAIL. They fail the original requirements concerning no harm and useful neural information. All five separate claims about history, order use and extra-task benefit also FAIL their required combinations of conditions. Some aggregate error reductions do not cancel failures elsewhere. If the rule requires every part to pass, a better average cannot rewrite it.\par

The negative report documents that these methods failed their rules. That result does not establish that either method is effective. It changes the next decision: neither candidate can be treated as the qualified predictor on which a teaching claim rests. The failures give the next experiment something specific to explain; an average improvement cannot make those conditions disappear. \href{research-evidence.json}{Programme computational evidence summary}.\par

\section{Next Experiments}

We are still short of a qualified predictor, before asking it to teach anyone. The failed comparison leaves Brain2Body's learning question open. Its measurements cover only two normalized finger-sensor groups, without calibrated joint angles, forces or a complete body state, and neither recurrent candidate has been admitted. We can propose a constrained controller by stating what it observes, which actions it may issue and what feedback it receives. If we give it the desired goal, we must report that supplied help when interpreting the result. Otherwise we could credit the predictor for an answer we handed to the controller.\par

The next comparison starts with the grasp itself. Before asking whether a model has read the right brain signal, we can ask whether learning movement patterns helps it predict another person's movement. The proposed development study uses qualified movement recordings from two people. One supplies source movement; the other supplies target training and development observations. It has not been run.\par

Give the model three different beginnings. First, let it practise on the target person's training movement. Second, let it practise on the source person's movement with the body features put in deliberately incorrect correspondence. Third, use the proposed correspondence between those source and target features. Keep the architecture, random starting conditions, number of fitting updates, later target-training examples and final-checkpoint rule matched. The wrong-mapping comparison asks whether a useful correspondence matters; the target-only comparison asks whether borrowing another person's movement helps at all. Neither comparison assumes transfer will succeed.\par

The proposed representation follows sensors at the wrist, index finger, thumb and object. Each position relative to the wrist has three coordinate components. The index, thumb and object positions therefore supply nine values; the distance between the two digit sensors adds one. Three components of wrist velocity complete the thirteen movement values. The wrist frame must be separately qualified, so its orientation is known. Wrist velocity describes the wrist's movement through the world, expressed in that frame. Its estimate must use only observations available at the current moment or earlier, never a later observation. That is what causal means here.\par

A characteristic sensor-distance scale separates size from shape. Fit it only on each person's training observations, then hold it fixed: development observations must not help choose it. This scale is not a measurement of bone length. Divide each position or digit-separation value by that length. For a velocity component, first multiply metres per second by the same one-tenth-second reference interval, then divide by the length in metres. The units cancel, leaving dimensionless features that put the two people's measurements on a common numerical scale. The reference interval does not change the recording's timing or establish that the bodies have identical dynamics. Physical errors would still need reporting in their qualified original units.\par

Each of the thirteen input features also carries a flag saying whether it was observed. That lets the model distinguish a missing value from a genuine observed zero. Future outcomes have their own observation flags for scoring: a missing target must not be counted as a measured zero, and a batch with no observed targets cannot supply the defined training loss. These future-target flags are not extra information for the prediction input. The units, sensor calibration, clocks and gaps remain questions to resolve before using the recordings.\par

The wider programme can compare both corresponding anatomical features and different features that perform roughly similar mechanical jobs. This two-person study tests one human sensor-landmark correspondence. It does not test whether different species have equivalent neural channels. Learning a movement constraint from another recording also does not turn that recording into a brain signal measured alongside the target person's action. A scan may help constrain anatomy where it fills a stated gap. A scan from elsewhere, or movement generated by a simulation, cannot provide the missing simultaneous brain measurement.\par

Keep recordings from the same acquisition, their overlapping pieces and their derivatives together when assigning roles, before making windows or calculating training statistics. Fit the size scales on training observations alone. Development observations can help assess this declared comparison; they cannot quietly become a final independent confirmation. Later neural adaptation would need real paired observations, the correct causal timing and its original input-specific controls. The protected recordings stay protected. Movement pretraining is a proposed intermediate test, not another successful result or evidence that a nervous system has learned a skill.\par

To test teaching, compare experienced, inexperienced and artificial trainers under matched conditions: the same task information, practice, assistance budgets and relevant burden. The interfaces would first need support from evidence. Include sham, delayed and unrelated-feedback controls with their designs and information access stated. Those comparisons could help test the contributions of useful information, timing and added input; their interpretation depends on the actual design. Better performance while the trainer is present would still leave the retained-learning question unanswered.\par

Then stop the trainer's information while keeping practice, feedback, task access and any necessary movement bridge constant. Measure performance later and on unfamiliar variations against matched controls, and record the help that remains. A gain that disappears after teaching stops counts against retained learning. A gain that survives needs its own limits established; it cannot by itself show tissue repair, a restored pathway or regeneration, which require separate biological measurements.\par

\section{Independent Verification}

Training software calculates how to change a model. Automatic differentiation computes derivatives, information describing how a change would affect a calculation. Adam is an update method that uses derivative information and stored state to adjust model parameters. Both training candidates rely on common Torch machinery for parts of these calculations.\par

Agreement between two routes using the same machinery cannot independently certify the shared machinery itself. This does not prove the machinery is wrong. It defines what the agreement has and has not checked. A stronger verification route needs independent calculations and complete state evidence appropriate to the actual training interface.\par

Reopening a saved model leaves a further question: can its complete state be recovered, and do subsequent updates match the saved training trajectory? Evidence for complete model and optimizer state, random-number-generator state, reload and the training trajectory remains unverified. A random-number generator supplies the numbers used for declared stochastic choices, so its state matters to that comparison too. Full independent NumPy calculation also remains unverified; NumPy is a separate numerical software library that would provide another route for checking the calculation.\par

Continuing training would require independent checkpoint and update evidence; the current record does not establish a verified continuation. No new fit, optimizer execution, reserved-test evaluation, accepted model or full-training GPU promotion is claimed here. All four reserved dates remain CLOSED. A GPU is a graphics processor that can also perform some computation; using one does not prove better scientific results or faster complete training. \href{research-evidence.json}{Current evidence and limits}.\par

\section{Human Finger-Intent Classification}

An earlier human study used recorded finger attempts. After adjustment using the first 50 labelled examples from each recording session, the selected classifier correctly identified 116 of 150 later trials. A classifier predicts a category; calibration here means adjustment using those early session labels. The support from calibration is part of the result, not something to omit from its description.\par

The later keyboard replay asks what happened when those predictions reached a simulated hand. Forty-one of 60 source-labelled right-finger trials produced the intended key contact. The controller executed all 54 supported commands that were actually issued. Thirteen presses were unintended or mistimed. These counts describe different questions and different denominators.\par

It helps to imagine a delivery service. It might deliver every parcel to the address it was given. If some addresses were wrong, that excellent delivery score would not mean every parcel reached the intended person. Similarly, 54 of 54 command executions does not erase wrong or mistimed commands before they reached the hand.\par

The recordings did not measure matching piano performance, spinal activity or peripheral nerve activity. The hand model was generic, and only four right-hand finger classes had supported key actions. The predictor waited for its recording window to complete; its earliest answer came one second after the cue plus computation time. That is not an instantaneous live readout.\par

These inspected later trials cannot remain fresh confirmation if future tuning uses them. Their scientific role changes when their answers help development. The earlier 116/150, 41/60 and 54/54 results must remain separate from the current macaque study and its failed benefit gates. \href{research-evidence.json}{Programme computational evidence summary}.\par

\section{Simulated Keyboard Learning}

A separate simulated hand policy completed 32 of 32 scored phrase cases in its constrained task. These were eight reserved simulated orders tested at two tempos and two poses, not 32 independent people or animals. It used a generic four-key hand whose wrist, thumb and sideways finger motion remained braced. It received requested keys, phase, duration and fixed smooth fingertip goals. Those are supplied task resources. Human neural clips shown alongside it did not drive the hand.\par

The record retains an interrupted original evaluator and a completed independent replay check; the tests were not reopened. It also discloses that soft simulated key limits can be exceeded in actual simulated travel. The display does not clamp a pose to hide that difference. A scored task result does not remove the need for physical and anatomical reference checks.\par

A current saved development replay has 336 frames with a 20 millisecond control interval. Its reported task results include 22 of 22 correct strikes, six of six chords, 22 of 22 timely releases, a worst timing error of 34 milliseconds and no extra strikes. A chord means keys pressed together. These figures describe the saved simulated example under its stated conditions.\par

They do not establish that a neural policy learned to play music, that a human acquired a skill or that a macaque's brain drove the animation. A policy is the rule or model choosing actions; here the interface gives it explicit goals. The software can be useful while that biological inference remains unsupported.\par

The separate human EEG, ESG, ENG and EMG clips also do not drive a combined macaque/human arm simulation. Actual joint and force ground truth and physical unit calibration remain gaps. A normalized finger-sensor target is not a simulated joint angle, and neither becomes measured neural tissue merely because it appears in the same three-dimensional scene. \href{research-evidence.json}{Programme computational evidence summary}.\par

\section{Failure Analysis}

Saved failures can tell us where to look next. In one reported index-finger comparison, a decoder with an extra neural-prediction task had more error than the simpler baseline. The comparison uses mean squared error: take each prediction's difference from its target, square it, then average those squared differences. Much of the excess was associated with an average offset; a smaller part came from variation around that average. An offset means predictions tended to sit too high or too low relative to the target. Locating that pattern gives us a question to investigate, not a diagnosis by itself.\par

That description does not uniquely establish the cause. Different trials, changed recordings, calibration differences or genuinely different states could contribute. If we use the validation answers to shift the predictions afterward, we have learned from those answers. We cannot then call the adjusted result an untouched test under the original rule.\par

Training-only summaries examine variation within trials, between trials and between dates. A frame is a time-indexed observation. Overlapping frames share information and are not independent new subjects. Counting many of them does not remove the need to consider dependence when interpreting uncertainty.\par

Named scalar arithmetic has limited independent checks, with a rejected skewness calculation preserved and a corrected calculation checked. Skewness describes asymmetry in a distribution of values. Correcting that arithmetic supports the stated summary, not a claim that a new decoder works or that neural drift has been causally established.\par

The proposed context branch would add a small model component using earlier neural information. It has not been fitted. Choosing or normalizing that component from the validation answers would make it another use of those answers, so any next fit needs admitted training-only material, appropriate controls and independent state/update checks. The proposal remains a candidate change to test, not evidence that the failed model has been fixed. \href{research-evidence.json}{Current evidence and limits}.\par

\section{Current Viewer and Source Identity}

A graphical user interface, or GUI, is the on-screen system people use to inspect and control software. For this programme, its scientific job is to make the evidence easier to understand. It should show measured data, predictions, simulation and missing values as different kinds of information.\par

Saved replay is valuable. A reader can pause, move to a later moment or inspect one frame. But the replay must retain its original source-time clock and full-run metrics. A slider at the start should not pretend that a recorded event happened at some new scientific time. A striking frame should not replace the complete task result.\par

Cache-only playback means loading an explicitly identified saved result without quietly running a new experiment. If the selected replay is absent, the interface should show that absence. If its source fails, it should clear the old replay and geometry. Otherwise a stale successful picture may be mistaken for a current result.\par

Natural-language goals can make the interface approachable. ``Press this sequence'' can become an explicit task specification with timing and permitted actions. But sending files to an AI chat is not the same as giving that system an execution environment. The interface must state whether it is reading context, predicting, controlling or simply displaying saved evidence.\par

\section{Application Performance}

The latest software candidate reduces its single HTML file from 113,298,857 to 34,295,597 bytes by moving opaque resources into the package. HTML is a source format for web pages; the bytes are file size. The moved resources still need to be delivered as part of the whole online directory.\par

This is a static packaging change. It does not demonstrate reduced startup time, lower memory use, fewer crashes or a GPU gain. Those require matched measurements of the complete useful application and its behaviour. A file-size number cannot supply an unmeasured performance result.\par

The candidate's browser verification was unavailable, and it remains unpromoted. A separately compiled original-source offline export is a different artifact from the online candidate. It cannot be used as proof that the candidate's whole interactive system works disconnected.\par

Offline also has a concrete limit: research views depending on the application programming interface, or API, are unavailable without that service. An API is a way software requests information or work from another component. Earlier bounded playback checks leave the current whole interface, disconnected
use, anatomy and function unverified. \href{research-evidence.json}{Current evidence and limits}.\par

\section{Anatomical Reference Tests}

The viewer can help us find a missed contact only if its own body has been checked. Every visual model needs a series of real-world references: one appealing image cannot establish all proportions, surfaces or relationships. Bones, muscles, joints, neural routes and contact should be checked where they are claimed. The parts and the assembled scene both need review. Otherwise a convincing display could hide an anatomical error while inviting us to diagnose the trainee.\par

The target is a tangible, lifelike three-dimensional model. Transparent layers can reveal a muscle over a bone. A cross-section can show a route that would otherwise be hidden. Paired views can clarify the front and back. Biological identification colours, useful close-ups and tilted orthogonal views should clarify the relationships. An orthogonal view looks along a defined right-angle axis; tilting the scene can retain that relationship while revealing depth. Illustrative signal paths should be distinguishable from actual tissue. Each choice should explain something real, rather than disguise uncertain anatomy with visual detail.\par

A template brain and an atlas are reference tools. They are not automatically registered to the recorded subject. Registration means relating measurements to the same individual's anatomy under a checked procedure. A region coloured on a template is not proof that we measured its neurons or know their function.\par

The source reconstruction matches the named template geometry and atlas values. Its optional region boundary still has documented mesh defects. These checks establish that particular correspondence, not full biological or visual acceptance. The retained historical conceptual figure explains an earlier framework. It neither proves current scientific acceptance nor reverses the present model failures.\par

\section{Time Alignment}

The proposed hand--arm task would use synchronized left and right tripod views. A motion-capture relay-composite is the programme's proposed method for aligning or stitching such views. Its purpose is to make a movement easier to observe when one viewpoint hides part of the body. It is not yet a demonstrated universal reconstruction method.\par

Imagine one camera sees the fingers while another sees the underside of the hand. The cameras need calibration and a checked relationship between their times. If they capture different moments, an assembled pose may never have existed. Occlusion, something blocking the view, can still leave important surfaces or contacts unknown.\par

LiDAR uses light to measure distance. A perpendicular view may help measure a particular surface, depending on the object and instrument. It does not guarantee access to every surface, reveal neural signals or solve registration by itself. The acquisition method should be compared with suitable independent references.\par

The joined view has to be checked against a series of independent real-world references, with its source clock, physical units and uncertainty preserved. Those checks tell us whether the acquisition method reconstructed the measured movement under the tested conditions. A better view can help us ask that question while leaving hidden tissue and missing observations unresolved.\par

\section{Selecting the Next Skill Test}

Four keys are a useful beginning because selection, contact, timing and release can be scored. They are a small part of playing an instrument. Musical skill also involves richer timing, coordination, posture, sound and interpretation. Passing a four-key task cannot honestly become ``piano mastery.''\par

Threading a flexible strand through a hole tests a different missing capability: shape, friction, contact and small corrections. A puzzle requires object-state tracking and coordinated grasping. Archery requires forces and control across more of the body. Breakdancing adds balance, whole-body movement and impacts. The current finger study does not contain all of these measurements.\par

The appropriate next task should expose a weakness that the model can physically express and the current evidence does not make trivial. We should reserve new objects, starting positions, action sequences and recordings before choosing the learner. A controller given the correct answer directly is a useful privileged comparison, but it is not a neural model that predicted that answer.\par

Simple alternatives must get a fair comparison. They need comparable task information, practice and evaluation. A complicated system should show why it is needed. Report complete-task success, extra actions, misses and delays; do not choose only the clip in which the movement looked convincing.\par

Maze comparisons can test orientation and memory. Right-angle and non-right-angle routes have different demands, but a harder-looking maze does not automatically mean ten times the difficulty. Difficulty needs a declared measure and comparable conditions. Biological learning times cannot be inferred from how quickly a computer finishes its training.\par

\section{Body and Species Differences}

A hand, a paw and a four-legged body have different proportions and available movements. A useful common description might be ``reach this target,'' ``maintain contact'' or ``release before the next step.'' That lets us compare task relationships without pretending the muscles or neural channels are identical.\par

Think of tailoring a pattern to different bodies. A shared pattern may be adaptable, but averaging two bodies does not guarantee a garment that fits either. The programme should preserve each organism's anatomy and test the adapted mapping against that organism's mechanics.\par

Physical size, computational detail and statistical grouping are separate. A tiny animal can have a highly detailed model. A large body can have a coarse one. Averaging errors over several dates can hide a failure on one date. The word ``micro'' cannot safely stand for all three at once.\par

Cross-species transfer is therefore a research question. Which relationship survived? Which required adjustment? Which failed? If no shared mapping works, a set of useful source-specific systems may still be an honest result. The current macaque hand mapping does not demonstrate biological teaching across species.\par

\section{Developmental Stages}

The native-development branch studies an organism developing with its own tissue. A body and its nervous system change as they mature. That may change recording access, what movements are possible and how a mapping should be tested. It does not make young tissue an automatic solution to learning, healing and interface stability.\par

Cook and colleagues documented postnatal development in the fat-tailed dunnart, including different developmental rates across body parts. These observations are species-specific. They do not establish the same timing in antechinus or prove a safe neonatal neural interface. \href{https://www.nature.com/articles/s42003-021-02506-2}{Cook et al., 2021}.\par

Armstrong and Ferguson studied skin-wound healing in young \emph{Monodelphis domestica}, an opossum. The transition in scarring is evidence about skin, age and that species. It does not demonstrate scar-free neural recording, spinal repair or lifelong stability of implanted electrodes. Dunnarts, antechinus and these opossums must not be silently treated as the same animal. \href{https://pubmed.ncbi.nlm.nih.gov/7750642/}{Armstrong and Ferguson, 1995}.\par

Artificial rearing and altered development would raise further questions requiring separate evidence. Neither has been demonstrated as a requirement for useful NeuroTranslation. Ordinary recordings and simulations can already test limited questions, so those questions can proceed while the developmental proposals remain unresolved.\par

\section{Proposed Twin Comparisons}

The proposed programme includes identical twins, non-identical twins and unrelated comparisons. Those categories offer different ways to investigate variation. They do not prove that two nervous systems use interchangeable channels or that a skill can be copied between them.\par

Even closely related trainees have different experiences, physiological states and electrode placements. To compare them, the experiment would need matched tasks, developmental stages, body measurements and acquisition conditions. A difference in performance might reflect practice or recording quality rather than the relationship being tested.\par

The twins model must be tangible. It should show where recording, feedback and body access would occur, with each component and the assembled scene reviewed against a series of real references. A picture of linked trainees explains a proposed comparison; it does not show that the experiment has happened. The measurements would determine whether a mapping actually transfers.\par

The useful question is which sources of variation can be controlled and which remain. If a mapping transfers only after substantial new calibration, that requirement is part of the result. Similar genetics does not supply the missing functional evidence.\par

\section{Humanized Models}

Revah and colleagues studied human cortical organoids transplanted into rats. Organoids are small laboratory-grown tissue structures representing some aspects of the brain's outer layer. The grafted tissue matured and integrated with host circuits under the study's conditions. That supports defined questions about tissue and host interaction. \href{https://www.nature.com/articles/s41586-022-05277-w}{Revah et al., 2022}.\par

It does not make the rat's whole nervous system human. It does not transplant an expert's skills or establish repair of a complete brain-to-body pathway. The tissue, host species, developmental setting and intervention remain essential to understanding the result.\par

The paper has a later correction concerning a duplicated image panel. That correction is retained in the citation record; it is not a retraction or evidence for stronger abilities. If the original figure is used, its corrected version matters. \href{https://www.nature.com/articles/s41586-026-10943-4}{Author correction, 2026}.\par

Brain2Body keeps native development, defined human-cell contributions, mixed-species tissue and graft interactions as distinct proposed branches. Each should ask what a particular contribution changes and how that change can be measured. Combining impressive labels does not establish a faithful human model.\par

\section{Disease, Compensation and Repair}

An injury may interrupt an output route even when the trainee can still make useful neural decisions. A compensatory bridge could help express those decisions as movement. Another condition may change sensory input, muscle function or learning itself. The same intervention cannot be assumed to solve all of them.\par

The condition determines the experiment. Identify the impaired function, what the interface receives and changes, and the physical outcome that would count as a benefit. A result in one condition cannot settle another when its cause or remaining pathways differ. The interface has to answer the question posed by that particular impairment.\par

Temporary assistance means useful help while it is supplied. Retained learning means a performance gain persists after the relevant teaching help ends. Tissue repair means a biological restoration claim, requiring its own biological evidence. A person moving better through a compensatory interface does not by itself demonstrate that damaged tissue regrew.\par

These distinctions matter because several outcomes can be valuable. Assistance may improve daily function without teaching a new independent skill. Retained learning may occur while a movement bridge remains necessary. Repair might have to be tested with different measurements entirely. The current computational work demonstrates none of these biological branches.\par

\section{Biological Trainer Tests}

Brain2Brain is the programme's label for specified information links between nervous systems. Pais-Vieira and colleagues connected recorded activity and stimulation to share task-relevant information between trained rats. The animals had learned task elements beforehand. Their bounded choices do not show an untrained animal instantly acquiring another's complex skills. \href{https://www.nature.com/articles/srep01319}{Pais-Vieira et al., 2013}.\par

A proposed trainer study would compare an experienced partner with an inexperienced partner. Both conditions would need comparable practice opportunities. If either partner gives the same benefit, expertise may not be the important ingredient. General guidance, additional exposure or the mere presence of feedback may explain the result.\par

Sham guidance would supply input without the intended useful information. Delayed guidance would test whether timing matters. Matched feedback unrelated to the trainee's current need would help separate the amount of input from its informative relationship to the task. A simple machine controller would provide another comparison.\par

The study would record what guidance was supplied and what the trainee actually did. Independent scoring and appropriate assignment of conditions would reduce the risk of choosing only favourable trials. These are proposed controls, not experiments already completed. The outcome could be no advantage, useful assistance or a retained benefit in a limited domain.\par

\section{Artificial Trainer Tests}

An artificial trainer is a computer model proposed to supply useful guidance. It may learn from measured examples or use a supported task model. Its existence does not mean it can teach a living nervous system. The interface and the information it receives are part of the claim.\par

Suppose one trainer sees only the trainee's current neural recording, while another is given the correct next action. Their scores cannot be compared as though they solved the same prediction problem. A desired goal, an expert trajectory, a correct label and measured sensory feedback are different resources. All should be disclosed.\par

A fair study would compare the artificial trainer with simpler alternatives and the biological trainer, matching task information and practice where the question requires it. It would measure the burden of receiving the guidance as well. A direct neural route adds complexity; a useful outcome must justify that complexity under the comparison, rather than the comparison assuming it will help.\par

An artificial trainer may help while remaining active. That would be assistance, which can be useful. The stronger learning claim asks whether the nervous system becomes more capable without trainer commands. We cannot get that answer from a computer policy succeeding in its own simulated task.\par

\section{Trainer-Withdrawal Tests}

This is the key experiment. Stop the trainer's commands, then test what remains. But hold any necessary compensatory movement bridge constant. If the trainee has an interrupted output pathway, removing the bridge could stop movement even if learning had occurred.\par

Imagine teaching someone to type through an interface that lets their intentions reach a hand. Disconnecting both the teaching help and the hand interface would mix two questions: did the trainee retain the skill, and could they still express movement? The test would no longer distinguish them.\par

The baseline and retention test need comparable access to movement, task conditions and feedback. Record any help that remains. Compare immediate performance, later performance and unfamiliar variations. Improved performance only while the trainer operates is assistance. Improved performance without trainer information needs further checks before calling it retained learning.\par

If the gain survives, report where it survives and the uncertainty. If it disappears, report that. Tissue repair is still another claim with different evidence. Holding the bridge constant does not prove repair; it makes the learning question interpretable.\par

\section{Networks and Skill Transfer}

BrainNet is a metaphor for specified connections between neural systems. It does not name an operating network through which people can buy biological expertise. The programme's motivation includes broad enablement and possible future skill services, but neither is a demonstrated product outcome.\par

Useful access, ownership of biological information, deployability and commercial services are different questions. A conceptual article cannot decide all of them from a simulated hand. Nor does a bold name establish that the proposal is unique, universally superior or the fastest possible route to learning.\par

To test a proposed skill gain, follow what the interface observes, the guidance it supplies, the performance that improves and the help that remains. An experiment must be able to reject the explanation. That inspectable research direction is what the current proposal offers; the evidence has not delivered the broader ability yet.\par

% R152 source proposal; provenance and exact prior text remain in the proposal manifest.
\section{From a reach to a testable model}
Start with a movement we can measure well. For the proposed snap showcase,
the arm begins extended with the palm down. The elbow and forearm swing and
turn; the wrist and hand turn upward during the snap, then move down to rest.
The model must account for that whole movement, including the preparation,
contact and return. Finger pairing, organism-specific anatomy and timing still
need a real-reference series; a human-looking animation cannot establish that
a particular monkey can perform the same gesture.

The cup reach supplies a simpler precision comparison: relate recorded signals
to measured movement before attempting the coordinated snap. Each task still
needs its own anatomy, contact, timing and independent comparison.

That comparison connects three interacting groups: task-linked brain
observations, named-muscle observations and a musculoskeletal simulation.
A learner would relate the observations to the mechanics; independent movement
and contact measurements would test its predictions. EEG records scalp
electrical activity and surface EMG records muscle electrical activity.
A separately recorded nerve or spinal signal needs its own declared channel.
The proposed apparatus uses two arm sleeves, scalp EEG and synchronized,
calibrated left and right tripod views. Vendor, channel, bandwidth and clock
qualification remain open.

Intent to action is the target. By understanding and measuring the intermediary,
we aim to make that mapping as close to one-to-one as possible within a declared
task. Specify what counts as the intended action, the participant and session,
state and context, permitted history, tolerated error and independent reference.
The test then asks whether observing the middle reduces ambiguity and improves
prediction or qualified control. It does not assume a unique hidden intention
or claim that the best possible mapping has already been achieved.

The proposed learner must explain the recorded brain and muscle activity,
where the hand moved and when. These are proposed comparisons: no paired trial,
trained model, data acquisition, fine-tuning, reinforcement-learning rollout or
accepted checkpoint is established for this EEG/EMG/movement loop. Its sources
and permissions, physical and computational behaviour, clinical interfaces,
reference comparisons and held-out tests still need independent assessment.

The completed LINK comparison supplies a warning about what would count as
success. Its 18 runs and 720 saved epochs tested a different question; both
model-admission gates and all five benefit claims failed. Neither those runs
nor a future successful prediction would show that a trainee received sensation
or produced an action. Recording, prediction, simulated movement, sensory
delivery and physical action each need evidence of their own. EEG and EMG
recorders alone cannot write sensation, and no sensory delivery or physical
actuation is established for this proposal.

\subsection{Form, function and action need a measured middle}
A structural map constrains what a system could do. A NeuroConnectome maps
specified neural connections; its organism, coverage, resolution and source
would have to be declared. Following the reach through that map also requires
the signals and the processes that turn those constraints into an outcome.
Signals must remain identified as observed or inferred. This is why the proposed
programme follows form to function to action: understanding and measuring the
middle could improve the intent-to-action mapping for the declared task.

The molecular comparison asks the same question about a different mechanism.
DNA to RNA and RNA to protein put a process between a starting constraint and its
downstream product. Regulatory context, RNA processing and cell state change
what that process produces; protein structure alone does not determine the
behaviour of a whole cell. AlphaGenome, RNA or cell-state models and AlphaFold
illustrate these different modelling jobs. AlphaTranscriptome names a desired
predictive capability here. Connecting the predictors would still leave system
dynamics and the environment to account for, with experimental outcomes needed
to test the combined prediction.
% Source-role pointers: alphagenome2026, scgpt, alphafold3; bounded owner checks.
\par\noindent \href{https://www.nature.com/articles/s41586-025-10014-0}{AlphaGenome regulatory prediction}~\cite{alphagenomeRetained};
\href{https://www.nature.com/articles/s41592-024-02201-0}{scGPT cell-state modelling}~\cite{scgptRetained};
\href{https://www.nature.com/articles/s41586-024-07487-w}{AlphaFold 3 structure prediction}~\cite{abramson2024}.
These examples support the comparison of modelling roles only. The complete
AlphaGenome and scGPT methods have not been independently assessed here, and
the examples provide no evidence of an integrated cell programmer.

For the reach, the corresponding biological question is what happens between
brain activity and the body's response. Brain2Spine and Spine2Body name the
proposed links to test; their composition is the Brain2Body aim. The spinal cord
integrates descending and sensory signals and contributes local coordination,
reflex regulation and adaptation. RNA processing is a different mechanism;
the shared question is how a context-dependent process changes an initial
constraint into an outcome. Supported spinal observations would test that
account. Where they are absent, internal states remain inferred or simulated.
Physical spinal work remains read-only.
\par\noindent \href{https://pubmed.ncbi.nlm.nih.gov/38603479/}{Spinal sensorimotor adaptation}~\cite{lavaud2024}
offers a constrained example of that processing. It validates neither our
proposed pathway nor unrestricted skill transfer. We draw on the abstract and
relevant author-manuscript passages; equivalence to the final published edition
and the complete methods and statistics remain unverified here.

Follow the action one step further and the question changes: what did it do
to the task environment? Thought2Action relates a declared intention to an
independently measured action. Action2World follows the action to its consequence
in that environment; Thought2World names the proposed composition. The earlier
Body2World label drew attention to the consequence outside the body. Body2Action
focuses on body mechanics and a specified task outcome, while Brain2Motion
focuses on the movement pathway. Each working label identifies a related scope
for a test; the labels establish neither interchangeable nor working controllers.

The engineering route would bring these observations and comparisons into a
learner: neural activity to musculoskeletal simulation to a programmable model,
toward NeuroTranslation. The proposed movement-focused learner is a large
movement model. It would perform the declared prediction or control. System 3
thinking names the proposed human--AI supervisory partnership around it: a
supervisory model would compare task strategies, diagnose where neural or
mechanical predictions miss independent measurements, and propose training or
calibration comparisons for the signal mapping and virtual body. The learner
would answer the task; the supervisor would help choose what to test next.
Neither replaces the biological spinal intermediary or a physical interface.
Final held-out outcomes must remain outside rewards and model selection.
\par\noindent Shaw and Nave's SSRN preprint,
\emph{Thinking---Fast, Slow, and Artificial}, and its
\href{https://executiveeducation.wharton.upenn.edu/thought-leadership/wharton-at-work/2026/05/thinking-fast-slow-and-artificially/}{institutional discussion}
provide conceptual context for this proposed human--AI role.~\cite{shawNaveSystem3Retained}
We use that framing only; the complete preprint and institutional article have
not been reviewed here. Neither establishes a third anatomical brain system or
validates the proposed learner or supervisor. Expert footage may supply movement
and tactical examples, with rights and camera uncertainty accounted for; it
supplies no measured expert EEG, EMG, muscle force or stimulation plan.

The outcome then supplies the next question for the learner and its supervisor.
Action changes what happens; its sensory consequences inform perception and
the next intention. Independently measured movement would anchor the action
part of that loop, while known stimuli, measured feedback and appropriate
recordings or reports would have to test the further sensory mappings. Movement
alone would not decode a unique sensation or perception. The return path must
be measured rather than drawn into existence. The full Interoceptomics aim
includes internal-body sensing and connectivity; external touch and
proprioceptive tasks contribute evidence but do not by themselves solve that
wider problem. A prediction, assisted execution and retained skill without
assistance answer different questions. Each wider claim must earn its own
independent comparison.

\subsection{Context connects. Grounding checks}
Within a reach, context has a concrete job: relate brain activity to muscle
activity in this task and at this time. Attention is one proposed way to
represent those relationships. The model compares a representation called a query with representations called keys, then uses the resulting weights to combine their values. Masks keep unavailable observations and future
samples out. The calculation weights information; it does not establish that
attention understands all context.
% Source registry IDs: context-native-01a12f62, context-attention2017-v7, context-mult2019, context-policy4
\par\noindent \href{https://arxiv.org/html/1706.03762v7}{Vaswani et al. (2017)}~\cite{vaswani2017}; \href{https://aclanthology.org/P19-1656/}{Tsai et al. (2019)}~\cite{tsai2019}. These methods motivate the computation; validation of the proposed paired system still requires the independent observations below.

A pattern can fit inside the model and still describe the wrong movement.
Grounding therefore means testing the declared prediction against qualified
observations under the same conditions. Independently measured movement would
check the outcome, while muscle activity would add an observation that
constrains the brain model. Even a matching reach cannot identify a unique
private experience. The proposed paired trial would have to show which of
these observations adds useful information.

\subsection{Muscle activity tests the brain model}
Record brain and named-muscle activity during the same reach. In the proposed
paired trial, EEG would sample scalp electrical activity and surface EMG would
sample activity associated with the named muscles. Paired training could use
the muscle observation to constrain their relationship. Then withhold the
muscle target and movement reference from the EEG-only predictor: can the
permitted brain observations predict a specified muscle-activity observable
or movement?

Pairing and calibration remain proposed. Shared cues and movement artefacts
could give that predictor a shortcut, so the comparison needs controls for
both. The muscle observations would cover only the named muscles, and EMG
supplies neither afferent nerve activity nor direct force. Sensory claims
still need known stimuli and measured feedback. Several neural states could
remain compatible with even a well-matched pair; the pair could narrow a
task-specific explanation without revealing the complete hidden brain.
Checking the movement itself requires the next independent observation.
% Source registry IDs: context-native-01a12f62, context-policy4, overhaul-brantley2018
\par\noindent This proposed pairing still needs supporting evidence. The candidate reference at \url{https://www.nature.com/articles/sdata2018133} has not yet been checked for exact support or reuse rights.

\subsection{Measured motion tests the virtual body}
Measure the same hand and arm from the aligned left and right tripod views.
The proposed motion-capture relay-composite would align the two point maps
in a common clock and coordinate frame, then join them into one composite
while retaining each point's origin. Before that joined map can check a
prediction, its calibration, alignment and capture precision would have to
be independently verified.

The movement map would constrain the external movement model: surface motion
and pose in the task's physical frame. To explain that motion mechanically,
the internal musculoskeletal model would describe specified bones, joints
and soft tissues. Organism, developmental stage and task would determine the
appropriate anatomical sources. X-ray projections could inform bone and
joint relationships; appropriate MRI references could inform specified soft
tissues. Registration, rights and source-specific limits would need independent
checks. No scan set is selected for this proposed comparison. Anatomy still
leaves activation, forces and neural function to test.

With a qualified composite, the musculoskeletal-simulation physics-engine's
joints, muscles, contacts and physical limits would face a measured comparison.
Fit the model on training data, then compare its predicted reach with
held-out measured movement. A force claim would require an additional force
measurement. The body's form limits possible motion; measured function would
test the mechanics and show where they need refinement.
% Source registry IDs: context-native-01a12f62, context-component-names-01a12f81, context-policy4, context-opensim2018, context-myosuite-docs, context-precision-01a12f77
\par\noindent \href{https://journals.plos.org/ploscompbiol/article?id=10.1371/journal.pcbi.1006223}{Seth et al. (2018)}~\cite{opensim2018}; \href{https://myosuite.readthedocs.io/en/latest/}{MyoSuite documentation}~\cite{myosuiteDocsUndated}. The simulation tools supply methods for that comparison; the proposed capture and body models still require independent validation.

A lifelike hand could expose a missed contact or a late reach. The comparison
still needs independent capture alignment, precision, mechanics and biological
fidelity; none has been demonstrated for this proposal. Captured geometry
cannot establish hidden neural state, internal forces or unique subjective
experience. It would check the declared movement prediction within those limits.

\subsection{Two pairs. One conditional learner}
The learner would have two connected jobs: relate EEG observations to muscle
activity, and relate measured movement to predicted mechanics. Large movement
model names the proposed family; deep neural network describes a possible
implementation. Large refers to intended scope, not a demonstrated parameter
count. Encoders would accept only observations permitted for the named test.
Two branches would summarize those relationships in a shared task
representation, with separate output parts for each direction's targets.
% Source registry IDs: context-native-01a12f62, context-policy4, context-attention2017-v7, context-mult2019
\par\noindent \href{https://arxiv.org/html/1706.03762v7}{Vaswani et al. (2017)}~\cite{vaswani2017}; \href{https://aclanthology.org/P19-1656/}{Tsai et al. (2019)}~\cite{tsai2019}. These computational methods offer precedents for a design, leaving the proposed paired learner and its performance to be tested.

Signal features belong to branches, and branches belong to the learner.
That hierarchy describes the computation, not two anatomical thinking systems.
A common clock, task and physical frame would connect the branches; independent
measurements would check their errors. Form constrains the prediction, and
checked function would guide refinement. Several hidden states could still
fit the same reach, so neither a shared representation nor conditional
predictions would establish an exact inverse nervous system. The architecture
has not been selected, implemented, trained or independently accepted.
To choose it, we would have to compare it with simpler baselines on
appropriately reserved observations.

\subsection{Learn the movement, then test what transfers}

A recording of a reach already contains a possible lesson for a computer: use the earlier movement to predict where the hand goes next. It need not contain EEG to pose that question. A separate proposed comparison would ask whether learning from one person's grasp-and-lift movement helps predict another person's movement. The WAY-EEG-GAL description identifies continuous recordings of lifting series as well as separate windowed trials and lift-summary tables.\citep{luciw2014way} The continuous movement could supply both the history and the later checking observation. It would not require matching AllLifts Run or trial rows, and it would create no measured neural pair between the two people. The original study of five complete weight series remains a different design.

To transfer anything useful, we must say what corresponds. Homologous structures share an evolutionary anatomical origin. Mechanically analogous structures can perform a comparable job without that shared origin: a robot gripper's opening and a human thumb--finger opening might both describe how an object is approached. That example would require a separate robotic study; the proposed two-person study would not test it. Each body keeps its own joints and possible directions of movement, positive-axis conventions, proportions, muscle--tendon roles, loads, inertia, contact and timing. Two similar-looking limbs need not generate the same forces or move at the same rate. Calculating a corresponding movement for another body is retargeting; the calculation remains derived, alongside separately identified tracking observations, estimates, predictions and simulations.

The proposed description of movement contains thirteen quantities. Three give the index sensor's position relative to the wrist, three do the same for the thumb and three for the object. One describes the distance between the two digit sensors; the final three describe wrist velocity using only information already available. Thirteen extra mask values say which features are actually available. Each person needs checked position units, sensor placement and a wrist orientation that turns positions into the intended frame without reflecting the body. After the checked conversion to metres, the middle, or median, wrist-to-index-sensor distance from that person's assigned training support supplies a positive length scale for position and velocity. This is a sensor-distance scale, not a measured bone length; development observations do not set it. Velocity is also multiplied by the same 0.1-second interval in both people. The thirteen prediction features then have no physical units. To recover velocity errors in metres per second, multiply by that person's length scale and divide by 0.1 seconds. Missing features contribute no target error; writing a masked zero does not make a missing observation measured. These features do not supply all finger joints, individual muscle forces or the person's inertia.

The model would use fifty movement frames at 100 observations per second to predict 100 milliseconds beyond the latest filtered input actually available. The source rate must first be checked as 500 observations per second. A proposed filter combining 51 present and earlier samples would precede keeping every fifth sample, producing 100 observations per second. Its linear-phase delay is 50 milliseconds. That delay stays in the account of when the information became available; moving a timestamp backward cannot make the predictor respond earlier. Filter coefficients, the frequencies it preserves, timestamp conventions, gaps and handling at recording edges still need checking. Only the named movement columns may enter. The source's kinematics matrix also contains force-plate signals, which cannot quietly become movement features.

Knowing which input was read is different from knowing whether two inputs are independent. Trustworthy published participant attribution, an unchanged exact archive entry and a checked capture can identify the participant record and recorded input without requiring one particular spelling of an ID inside the file. They do not authenticate a biological person or prove that two files came from different original acquisitions. A longer recording and its trimmed copy could have different names and file fingerprints while sharing the same movement. Shared samples, trims, derived versions and overlaps must therefore remain together as one acquisition group before filtering, cutting windows, scaling or assigning roles. The complete public non-ST roster and these relationships still need establishing. All 294 rows of the earlier AllLifts table remain unavailable for training, evaluation and assigning trials or acquisition groups. ST, competition storage and reserved LINK recordings remain closed.

Qualified P1 groups would provide the earlier movement learning. P2 groups would be divided between target training and development checking. At least two independently established P2 acquisition groups are needed. If there are $G\geq2$, a fixed hash calculation ranks their established original-parent identifiers; the first $\lfloor3G/5\rfloor$ become training groups and the rest become development groups. The floor brackets mean rounding down, so the nominal 60/40 division is explicit even when the number of groups is small. The exact identifiers, how they are written as bytes and the complete roster must be settled before ranking or seeing outcomes. The order is not chronology and need not balance lifting conditions. If uncertain ancestry forces all P2 recordings into one group, the comparison stops. Cutting that group into windows cannot create independent training and checking acquisitions. No actual group or movement segment has been assigned by this proposed identity method.

The three training arms keep the target task and model the same. B0 first learns from correctly mapped P2 training movement. B1 and B2 first learn from exactly the same P1 windows in the same order. B2 uses the proposed correct correspondence. B1 deliberately uses the thumb sensor for the index role, the object sensor for the thumb role and the index sensor for the object role, leaving the wrist anchor fixed. Both use the same correct source calibration and scale before this change, then the same correct target adapter and ordered target-training batches. Sampling gives each original acquisition group balanced treatment and records repeated as well as unique windows. The learner might undo the false correspondence; changing roles can also change which features are available. Equal numbers of updates do not give the arms equal diversity of source movements. Those limits must stay in the explanation of any difference, rather than prompting a more convenient control after results are known.

The proposed network is a single gated recurrent unit, or GRU, layer: a calculation that updates a numerical memory as it reads the sequence. It has a memory width of 32, accepts the thirteen movement features and thirteen masks as 26 inputs, and predicts thirteen outputs. That software memory is not measured biological memory. Each arm uses the same initialization for a given fixed random seed, with seeds 17, 31 and 73. Every arm/seed receives 200 earlier-learning updates and 100 target-training updates, using batches of 16 examples. Adam, the proposed weight-update rule, has learning rate 0.001, beta settings 0.9 and 0.999, epsilon $10^{-8}$ and zero weight decay. There is no clipping, changing learning-rate schedule or early stopping; the final update supplies the checkpoint. During training, squared prediction errors are averaged over observed target features; a batch with no observed targets must stop. Across three arms and three seeds this would be 2,700 updates. None has been performed for this comparison. No external checkpoint is selected, and the numerical updates, filter, precision and actual computing limits still need independent checks.

For each development acquisition group and seed, the proposed score adds the squared errors for the observed target features, divides by how many such targets were observed, then takes the square root. It then averages those group scores equally, so a long recording does not automatically outweigh a short one. A comparison with no observed targets has no valid denominator and must stop. The report must give feature counts, missing observations, errors converted back to physical units, uncertainty and a persistence forecast---continuing the latest available movement value---using the same information endpoint. Every worse group or feature remains visible. B2 would need lower normalized error than both B0 and B1 in each fixed seed, with no applicable original requirement failing, to justify another investigation. Two people and development checking do not establish performance across a population, untouched final confirmation or biological teaching.

Larger movement collections could supply more examples after their exact files, permissions, relevant movements and original-recording overlap are checked. \href{https://www.addbiomechanics.org/download_data.html}{AddBiomechanics} distinguishes recorded markers and ground reactions from fitted joint movements and torques estimated by inverse dynamics---calculating forces or torques consistent with motion and a body model. Those estimates are not new direct force measurements. The selected \href{https://amass.is.tue.mpg.de/license.html}{AMASS} asset would need its research-use, non-commercial and redistribution terms checked. The same original recording in two collections stays one source, not independent training and checking data. LaBraM's earlier training inventory flags the WAY source family, but exact checkpoint, participant and trial exposure remains unknown. A later decision to reserve a recording cannot remove its possible influence from already trained weights. No such checkpoint is selected for this comparison.

Scans have a more specific job than making the body look complete. Before adding X-ray/CT or MRI, check which anatomical sources the chosen musculoskeletal model already uses. A new scan is useful when it fills a stated gap or provides a genuinely independent check. Its person or species, scale, registration, permissions and uncertainty still need checking. Registration establishes the correspondence between that image and the body being studied. Another person's scan remains an external reference; a generic body model does not secretly measure this participant's geometry, muscle activation or forces. No scan set is selected here.

The actual eligible recordings, acquisition groups, units, frames, clocks and calibration are still unresolved for this continuous comparison. It reports no new fit or successful transfer. A movement description that proved useful could later be tested alongside genuinely paired neural and body observations, using only the inputs allowed for that prediction. The transferred or calculated target would not become a simultaneous measurement. The useful next answer is where the movement forecast improves, where it worsens and which measurement or assumption explains the difference. Neural equivalence, causal physiology, delivered sensation and a skill that remains after guidance ends would each need a further experiment.

\subsection{Pretrain. Pair. Verify. Refine}
Those comparisons start with the observations we may use. Data rights,
calibration and held-out subjects, sessions and tasks determine which can
enter the study; resolve them before processing data or choosing a model.
Open-source code does not grant blanket permission to use datasets,
checkpoints or clinical interfaces. Qualified data could first train
representations for each recording type. LaBraM is one EEG encoder candidate,
not a complete brain-body writer. Synchronized paired observations could
then train the declared mapping from permitted inputs.

Paired prediction would give later refinement a baseline to beat. Before
optional reinforcement learning, post-training or task-specific fine-tuning,
the physical predictions and reward need checking against independent
outcomes; a reward the computer can calculate supplies no biological proof
by itself. Offline reinforcement learning needs qualified coverage of actions,
transitions and rewards: labels alone cannot tell us what an untried action
would do. Backpropagation can update differentiable neural components, but
it cannot automatically pass through every physics engine; universal
differentiability is not established. Simpler baselines and frozen final
tests remain through every phase.
% Source registry IDs: context-native-01a12f62, context-policy4, context-labram2405, context-cql2020-abstract, context-opensim2018, context-myosuite-docs
\par\noindent \href{https://arxiv.org/html/2405.18765v1}{Jiang, Zhao and Lu (2024)}~\cite{labram2024}; \href{https://arxiv.org/abs/2006.04779v3}{Kumar et al. (2020)}~\cite{cql2020}; \href{https://journals.plos.org/ploscompbiol/article?id=10.1371/journal.pcbi.1006223}{Seth et al. (2018)}~\cite{opensim2018}; \href{https://myosuite.readthedocs.io/en/latest/}{MyoSuite documentation}~\cite{myosuiteDocsUndated}. These sources offer candidate methods for the sequence; validation of the proposed paired system still requires the experiment below.

This proposed paired training sequence has not been carried out.

\subsection{Respond quickly. Improve carefully}
A frozen accepted model would provide the proposed response: predict from
permitted past and present observations, then prepare encoded outputs for a
qualified interface. Inspecting the errors would be a separate job. That path
would calibrate and propose a candidate update, to be independently evaluated
with a restorable checkpoint before it could replace the running model.
Updating an estimate of the current state therefore differs from changing
learned weights; examining an error would not silently change the predictor.

The candidate update would need more than a good-looking new score. Its
independent evaluation, a saved version that can be restored, evidence that
restoration works and the applicable authority would all precede replacement
of the running model. Until then the accepted weights stay fixed. This keeps
an attempt to improve tomorrow's model from quietly changing the one being
tested today.

The fast-and-slow comparison is a Kahneman-like analogy for this division
of work. It identifies no pair of biological systems and assigns no inherent
thinking speed to EEG or spatial data. Whether the frozen response is quick
enough would depend on the complete path from acquisition through filtering,
alignment, inference, output, feedback and rendering, including delay tails
and jitter. That end-to-end latency has not been measured. No accepted model,
real-time demonstration or running system that adapts under these rules is
established here, so the proposed separation supplies no latency or accuracy
guarantee.
% Source registry IDs: context-native-01a12f62, context-policy4

The complete timing test would include preparing the output and any later
sensory conversion or actuation, rather than stopping the clock when the
prediction leaves the model. Delay tails describe the unusually slow cases;
jitter describes variation in delay. Missed deadlines, uncertainty and the
effect of other healthy workloads would also need reporting. An average
inference time could hide the late response that matters in a catch.

\subsection{Write sensation. Enact action}
What reaches the trainee when the computer has made its prediction? A sensory
transducer would have to deliver a sensory output; a controller and actuator
would have to produce action in a virtual, robotic or later qualified biological
body. Input and output describe the computer boundary. Afferent and efferent
describe physiological direction relative to the central nervous system.

A display or haptic interface supplies sensory feedback, which differs from direct neural writing. Direct cortical writing could have a sensory role without being a peripheral afferent signal, just as a robot command is not an anatomical efferent nerve. EEG and EMG recorders do not write, and predicting an EEG-like pattern does not deliver sensation. Each proposed output needs an independently qualified interface, permission for its use and a measurement of what it delivers. No such physical interface, stimulation, acquisition or actuation is established by this proposal.

No methods evidence from the inaccessible clinical source supports this
proposed interface. Implanted precedents cannot establish that EEG or EMG recorders are equivalent to devices that deliver sensation or produce an action. Those physical additions remain unpriced and outside the acquisition and virtual-modelling scope.
% Source registry IDs: context-write-native-01a12f66, context-native-01a12f62, context-policy4

\subsection{Close the loop with evidence}
Return to the cup with two independent checks: withheld body observations
for the brain-to-body prediction, and the recorded neural target for the
conditional estimate from the body. Agreement between the models or a high
reward could still accompany a wrong answer. Neither establishes a unique
inverse brain or someone's private experience. A later sensory or action
interface would need a further measurement of what it delivered, under which
conditions, with what error and delay.

Approximation gets us towards plausibility and virtualization gets us to
precision. Independent measurements, repeatability, causal checks and stated
error bounds must earn that precision; the proposal supplies neither perfect
biology nor an automatic result. The original held-out and real-reference
controls would test where each prediction works and where it fails. The
retained-learning controls ask a further question: what remains when useful
guidance stops, with any necessary movement bridge kept constant?

Physical sensory writers, transducers and actuators remain unpriced additions
outside the acquisition and virtual-modelling costs. Those costs cannot tell us
the price of delivering sensation or action. No completed loop, training or
physical-writing result is claimed here. \emph{Catch it again when the help is
gone} develops that task and retention comparison. \emph{Build a brain--body
system we can test} asks what those results would have to establish before the
prototype could grow. The mechanism, task and prototype are stages of the same
proposed experiment.

\subsection{What the attention calculation actually does}

The proposed learner needs a way to relate permitted observations. Consider
one moment in a reach. A query represents what the model is looking for;
keys represent the observations it could compare with that query; values
carry information from those observations. The model learns these
representations. Their names describe a calculation, not a person inside the
machine who understands the task.

The technical edition writes that calculation as
\[\operatorname{Attention}(Q,K,V)=\operatorname{softmax}\!\left(QK^{\mathsf{T}}/\sqrt{d_k}+M\right)V.\]
Here Q holds the queries, K the keys and V the values. The superscript on K
means transpose: rearrange its rows and columns so the comparisons can be
calculated. Multiplying Q by that rearranged K gives scores for the
query--key matches. The common width of each query and key is $d_k$; dividing
by its square root scales those scores. The softmax operation turns the
allowed scores into weights that sum to one for a query with usable admitted
keys. Multiplying by V then combines the values using those weights.

M is the admitted mask. It controls which observations may contribute, so
an unavailable muscle target or a later sample cannot quietly provide the
answer. A mask is part of the declared test, not a decorative option after
the model has seen the data. Channels, recording type, task and time must all
retain their meanings through the calculation. Even a correctly calculated
weighted combination could predict the wrong reach: the independent
measurement still has to check it. This is a proposed method; no weights,
EEG encoder or gradient route have been selected or executed here.

\subsection{What each direction is allowed to know}

Use the same short labels as the technical method: B means EEG, E means
EMG, Y means independently measured movement, C means declared task,
stimulus and context, and P means the physics model with its identified
source. The letters keep track of information. They do not make the
quantities interchangeable.

An EEG-only test would use permitted past and present B, together with
declared context, to predict a specified E or Y target. Training may have
used qualified paired examples. At prediction time, however, the withheld
muscle and movement answers and future samples must stay outside the
predictor. Alignment, filtering and preparation of the data must respect
that separation too. Supplying a target through a less obvious route would
still supply the answer.

A body-only test would use permitted E, Y or both to estimate specified
recorded neural features. Its answer is conditional: given these observations
in this task, what recorded features are plausible, and with what uncertainty?
Several hidden neural states could fit the same movement. The reverse
estimate therefore needs its own recorded target; it cannot recover a unique
inverse brain or deliver the sensation it predicts.

A hybrid test would use additional permitted recording types and need its
own named route and evaluation. Its result could not be presented as EEG-only.
Inputs calculated by the simulator would remain simulated. They would not
replace independent observations or inherit biological validity from a
lifelike hand.

Two models could also learn to agree with one another. That agreement is
called cycle consistency when following one mapping and its reverse brings
the result back toward its starting point. Both mappings could still be
wrong about the measured body. Shared cues, movement artifacts, subject
identity, overlapping sessions and information leaked during preprocessing
need the original controls. Final held-out answers must remain outside
selection of the model, physics and reward.

\subsection{How the proposed training would be checked}

First decide which observations may enter the study. Establish the rights,
people and tasks, physical units, calibration, uncertainty and shared clocks.
Freeze the relevant participant, session and task splits, then check whether
earlier pretraining already included material that would compromise those
splits. Those decisions precede preprocessing and model selection.

Qualified recordings could then train a representation for each recording
type. In the candidate LaBraM route, a neural tokenizer turns signal segments
into codes; masked-code prediction trains a model to estimate hidden codes;
later fine-tuning adapts it to a declared task. These are computer-learning
steps. LaBraM is a candidate EEG encoder, not a complete brain--body writer.
Permission to use code, checkpoints, datasets and individual recordings
requires separate qualification.

Next, paired supervised learning would use named measured targets and
explicit masks. Supervised means the training examples provide the answers
against which predictions are compared. Simpler predictors, cue-only
predictors and artifact controls would help test whether the proposed neural
route contributes useful information. Raw recordings, processed features and
model-generated values would remain distinguishable.

Combining errors requires another decision. An illustrative training
objective could include errors in EMG, movement and recorded neural targets,
applicable physical constraints and optional consistency terms. Before
choosing their weights, each term needs its target definition, units and
measurement uncertainty. Otherwise changing a unit or weight could change
which error the optimizer mainly tries to reduce. Consistency asks whether
the outputs fit together; an independent measurement asks whether they fit
what happened.

Optional reinforcement learning and post-training would come only after
checking the physical predictions and reward against independent outcomes.
Success on the task, smaller calibrated movement error and valid physical
states are possible reward ingredients, not results already achieved.
Loss and reward weights would be declared from training and development
evidence. Offline reinforcement learning needs observations that support
the relevant actions, resulting changes and rewards. A labelled trajectory
alone cannot tell us what an unsupported, untried action would do.
Conservative Q-Learning is one candidate method for this stage; its name
provides no safety or performance guarantee here.

Task-specific fine-tuning would still face frozen final tests in both
directions, with uncertainty, baselines and every original threshold.
Failure would remain a result. An accepted version could then run with its
weights fixed; a later update would need separate independent acceptance,
restoration proof and applicable authority before replacing it.

Backpropagation calculates how to change differentiable parts of the
computer model. It is different from biological feedback. A physics engine
does not automatically supply the required derivatives: the actual
OpenSim/MyoSuite simulation and optimization route, or a separately validated
substitute, would need checking. No such training, reinforcement-learning
rollout or task checkpoint is created or certified here. The procedure makes
the proposed comparison inspectable while leaving its outcome open.
% Source registry IDs: context-native-01a12f62, context-write-native-01a12f66, context-precision-01a12f77, context-policy4, context-spec15, context-glossary13, b2b-circular-checklist-original

% Proposed full-depth accessible conceptual-article insertion; no integration.
% Inherit the selected master's sans-serif text and mathematics setup.
\section{A larger model still owes us a learning experiment}

A computer could become better at predicting where a hand will move while the person learns nothing new. That is why the proposed Large Movement Model needs two kinds of test. One asks whether broader training improves a specified prediction. The other asks whether using that prediction helps a trainee, and whether anything learned remains after the teaching commands stop. A better answer to the first question does not settle the second.

The programme's new working hypothesis is that increasing both the amount and the variety of qualified neural, muscle and movement data could help a model cope with new situations. Amount and variety need separate comparisons. Neither has produced a scaling result here, and the proposed body-wide world model has not been trained or validated.

\subsection{Four training sources, five component roles}

The broader proposal draws on at least four training-source groups. Neural observations form one group. Muscle and other myographic records form another; the actual kind of muscle measurement must always be named. Measured motion and displacement form a third. Examples generated by physics-based motion or musculoskeletal simulation form a fourth. A complete collection of each has not been acquired here, and the list leaves room for other useful information.

Suitable neural recordings could come from the brain, spine, peripheral nerves and relevant visceral sources. Their sites, instruments and task roles matter. The present emphasis is motor--sensory. The programme uses Myoception for sensing muscle-driven movement and the feedback that comes with it, including body-position and movement feedback, touch and other task-relevant sensations. This proposed programme use of the name establishes neither a new external scientific standard nor a simultaneous recording of the whole body.

Touch, body-position and movement feedback (proprioception), contact, force and task context each need the right observation. Knowing which tactile stimulus was delivered differs from recording a sensory nerve response or asking what the person felt. A contact sensor's event differs from a calibrated force measurement. Tracking a limb can help estimate its configuration without measuring how that position feels to the subject. EMG records muscle electrical activity; it does not directly record peripheral nerve or sensory activity. Measurements, reports, estimates and simulated values must keep those identities.

The fourth group makes the virtual body a source of training examples as well as a constraint on a forecast. A suitable simulation could generate trajectories, changes between states, calculated muscle--tendon and contact forces, or movement patterns used as a starting point for learning. Record which simulator and body-model version produced them, its parameters, numerical solver, controller and original sources. A controller requesting joint positions and a controller driving muscles do different jobs; the first must never be described as the second.

If real movement is transferred to a different body model, preserve the original recording and the transformation. The internal states or forces calculated from it still belong to the simulation. A generated force is not an independent force measurement. The generator needs calibration against suitable measured support, with uncertainty stated and adjustable settings fitted from training support rather than independent tests.

A proposed test would keep the real training data fixed and add simulation examples in one arm. The permitted information, computing work and original decision rules would be matched, then held-out real observations would check whether the added examples helped. Record how often each arm sees real and simulated examples and how many updates it makes; the same real recordings can still be presented different numbers of times. Count the amount and variety of simulation separately from real acquisitions. Those real tests cannot also be used to calibrate the generator, tune its dynamics, choose rewards or select the model facing them.

The simulation may pass on a useful relationship or its own mistakes in contacts, dynamics and correspondence between bodies. The comparison must be able to expose both. Generated examples cannot independently verify the biological assumptions that produced them. The methods, implementation and use conditions of any proposed generator remain to be checked; no frontier implementation or dataset is selected here.

The four source groups and the programme's five component roles answer different questions. Brain signals, muscle signals, measured movement, a virtual body and a learning model say what the parts do. The four groups say where training examples could come from. The virtual body can supply both examples and physical constraints, and the learning model can use several groups. Additional useful observations remain possible. None of these resources supplies a retained-learning result before the trainee has been tested.

\subsection{Count more examples without confusing them with more experience}

One long recording can supply thousands of short windows. Showing those windows repeatedly gives the model more training presentations, but no additional person, animal or original recording. The proposed study would keep track of those differences: original acquisition groups, subjects, recording time, trials, available target values, distinct windows and repetitions. Copies and overlapping extracts from the same acquisition would stay together when data is divided for learning and testing.

Variety also needs to say what varies. Another person, another recording day, another movement and another sensor each change the problem differently. Body region, organism, apparatus, environment and measurement quality matter too. Counting each collection as one dataset would hide these differences. The study would describe them rather than assume that a larger dataset count means a more informative collection.

An amount comparison would add eligible training support from a declared source distribution. A diversity comparison would keep a specified amount comparable while changing named aspects of its range. Crossing those comparisons could test whether more data helps differently when the sources are broader. The model's design, permitted information, prediction target, preprocessing and selection rules would stay fixed. The study would report repeated exposure, updates to the model and computing work actually used. Equal numbers of updates need not mean equal cost or equal information.

If broader sources also have different measurement quality or easier tasks, that is a combined change. Its outcome cannot be assigned to diversity alone. Exact recordings, budgets, random starting conditions, controls, uncertainty methods and decision rules would have to be fixed before the comparison. This proposed expansion does not change the already adopted continuous WAY-EEG-GAL development experiment or open any original protected test.

\subsection{Let the curve be wrong}

One possible relationship between amount and error is
\[
E(N)=E_{\infty}+A(N/N_0)^{-\alpha}.
\]
$N$ is the chosen training amount and must be greater than zero. $N_0$ is a fixed positive reference amount in the same units, so their ratio is a pure number. $E(N)$ is a nonnegative error on the stated held-out population. In this proposed decreasing family, the floor $E_{\infty}$ and the changing part's size $A$ must each be zero or positive, and the exponent $\alpha$ must be positive. The floor and $A$ have the error's units; the exponent has none. Positive $A$ gives a strictly falling curve as training amount increases. Zero $A$ gives a flat line even when the exponent is positive: the exponent alone does not settle the direction. If the error measure also has an upper limit, predictions must respect it throughout the stated fitting and prediction range. A fitted floor would be a parameter of this curve, rather than an established irreducible biological limit.

No exponent has been fitted here. The measured error is free to stay flat, plateau, rise, reverse or change differently across tasks. A constrained decreasing curve is one candidate explanation; it cannot rule negative transfer out of the results. Those alternatives need to be taken seriously before the results arrive. The study would set its fitting range, competing curve shapes, uncertainty calculation and checks on withheld training-amount levels in advance. A fitted curve must predict data levels left out of fitting. Many overlapping windows from one recording do not become independent repetitions because there are many of them.

Errors also keep their original meaning. An error in a processed muscle signal is different from a camera-pixel error or a joint-angle error. They cannot be mixed into one percentage of biological accuracy without a justified, independently checked scoring rule. Report each measure's units, observation counts, missing values, domains and uncertainty. A bigger computational model is another change again; it cannot be quietly bundled into a claim about more training data.

To test generalization, name what is new. An unseen day, an unseen person, an unseen movement, a new dataset and another species require different comparisons. Each proposed transfer would face learning from target data alone, an appropriate deliberately wrong correspondence and simpler predictors, along with the original tests for shuffled associations, shared cues, recording artefacts and future-information leakage. Poorer performance on particular outputs or domains stays visible. A better average cannot undo an original rule that no evaluated output may be harmed, nor erase exposure during earlier training.

\subsection{Keep the missing parts of a recording missing}

FALCON M1-A offers one candidate connection: motor-cortex activity recorded together with processed muscle electrical signals from electrodes within the muscles of one rhesus macaque. This is intramuscular EMG, which differs from a surface recording over the skin. The recording location and the later signal processing are separate facts. The reviewed release does not establish synchronized motion-capture trajectories. Its muscle targets use processing that includes later information, which limits what an offline result could say about a purely causal real-time prediction. Its small minival subset comes from calibration data and is used to check submission format. Those reused examples do not supply independent validation. \href{https://doi.org/10.48324/dandi.000941/0.241029.1405}{FALCON M1-A release}; \href{https://papers.nips.cc/paper_files/paper/2024/file/8c2e6bb15be1894b8fb4e0f9bcad1739-Paper-Datasets_and_Benchmarks_Track.pdf}{FALCON methods}.

AJILE12 offers another connection: intracranial recordings and upper-body positions estimated from video in twelve clinical participants. These positions are two-dimensional camera pixels. The reviewed source does not establish paired EMG or instrumented three-dimensional motion capture. Its pose estimator is adapted using each participant's labelled video. Whether a particular day later reserved for testing was used to train that estimator remains unverified, as does a numerical bound on alignment error. Labels may therefore be participant-adapted even when the neural prediction model has never trained on that participant. \href{https://doi.org/10.1038/s41597-022-01280-y}{Peterson et al., 2022}; \href{https://doi.org/10.48324/dandi.000055/0.220127.0436}{AJILE12 release}.

The two sources could contribute useful, different training relationships after qualification. They never become the same event. A macaque's brain--muscle recording and a person's brain--pose recording do not supply an observed brain--muscle--movement triple. A model might learn patterns separately, then try an explicit transfer to another source. Actual target measurements would have to test whether that correspondence works.

The release descriptions identify possible source relationships. They do not establish downloaded recordings or a trained model. Exact versions, use conditions, file integrity, original acquisition relationships, timing, units, earlier exposure and the independence of the chosen partitions still need checking before training. Mapping motion between bodies also needs declared geometry, axes, proportions, load, contact, timing and uncertainty. Bodies may share useful mechanical relationships without sharing identical neural codes.

\subsection{Make the inside explain the outside}

When fingers close around an object, the visible movement depends on what happens within the body model and at the contact. The proposed world model would connect specified joint and muscle mechanics to supported movement, objects and environmental conditions. Its interior initially means those named mechanical states. It does not mean that every organ, nerve or sensation has been reconstructed.

NeuroTranslate would map named observations to specified targets. MoCap ReComp would supply a separately qualified movement composite that retains where each observation came from. An ordinary set of estimated video poses does not become that proposed instrument by receiving its name. MyoSim Replay would calculate the declared body's mechanical response. Their proposed combination is the Large Movement Model. Bringing neural representations from different recording types into one model would not establish matching individual cells across organisms or assemble a measured whole nervous system.

A proposed internal state could describe joint arrangement, generalized velocity, muscle-model activation and the environment. Rotational positions and velocities may use radians and radians per second; sliding coordinates require length units. A transition model would use permitted earlier observations, context and declared controls to forecast possible later states and their uncertainty. Every instrument would have its own observation relationship. EMG supplies muscle electrical activity, not direct peripheral nerve activity, and a processed EMG amplitude does not automatically measure activation or force. A tracked position is different from an internal force calculated by the simulator. Missing observations need an availability mask, and numerical filling values must never acquire the status of measured zeros. Predicting recorded neural features from body observations can leave several hidden explanations possible.

The physics needs more than a believable picture. Predicted movement must be checked against independent measured motion in matching time, coordinates and units. A claim about internal force needs suitable force evidence as well: a similar path can be produced by different combinations of muscle forces. Joint limits, loads, contacts, gravity, friction, energy accounting and sensitivity to numerical settings need their own checks. OpenSim offers established methods for specifying musculoskeletal models; it does not validate this proposed assembly. \href{https://journals.plos.org/ploscompbiol/article?id=10.1371/journal.pcbi.1006223}{Seth et al., 2018}.

% R4 full-depth accessible conceptual-article insertion; no integration.
% Inherit the selected master's sans-serif text, code and mathematics fonts.
\subsection{One muscle pulse, two simulated trajectories}

A computer hand can respond to a muscle command before any person has learned a skill. This test asks about that response. We compared two trajectories of \texttt{myohand\_r}, an installed mechanism representing a human right hand attached to a fixed passive torso. Both started in the same recorded state; one received a short muscle excitation. The comparison tests whether the specified command changes the state the software computes. It contains no brain decoding, measured biological response or evidence of retained learning. The software was \texttt{myo\_sim} 0.2.3 at commit \texttt{93b0ca8f4ec90c9899ee7f05fee561e9911da91b}, MuJoCo 3.14.0 and NumPy 2.5.3.\cite{myohandMyoSimPinned2026,myohandMuJoCo3140,mujoco2012,myohandNumericalRecord20261105}

\subsubsection{Let the control move too}

The compiled hand has 23 scalar hinge coordinates, meaning 23 angles describing its movable joints, measured in radians. It has 39 muscle--tendon actuators that calculate muscle-driven pulls, with 39 internal activation states and zero sensors. These are parts of a computer mechanism. Its activation values are not measured muscle electrical signals, and its forces are not readings from a force instrument.

The two branches began with matching complete recorded initial states, including angles, velocities and muscle-model activation. One received no additional excitation. The other received dimensionless excitation 0.2 at the actuator named \texttt{FDS2\_r}; all other control values were zero. Dimensionless means that 0.2 uses the model's input scale, rather than volts, newtons or an EMG amplitude. It drives a muscle model instead of requesting a joint position directly.

Both branches used the native Euler setting, the simulator's specified rule for advancing its state. Each advanced 100 times by $0.002\,\mathrm{s}$, reaching nominal $0.200\,\mathrm{s}$. Starting the count at step zero, the pulse ran at steps 10 through 59. Its pre-step interval was $0.020\leq t<0.120\,\mathrm{s}$. Gravity and passive mechanics could move either hand. The unpulsed branch therefore supplies a moving comparison, rather than a requirement that the hand freeze when the added command is absent.

The state-dependent quantities were recalculated with a forward evaluation at the start and end of each branch. The saved record contains 200 explicit integration steps and four explicit forward evaluations. The executed construction path also implies two compilations and one internal forward evaluation while the mechanism was being built. These are counts of identified calls, not a tally of everything inside the native solver.

The construction settings need care. The optional \texttt{inertia\_floor} argument stayed at \texttt{None}, yet the native builder changed each diagonal rotational-inertia component of four fixed scaffold bodies from $1$ to $10^{-4}\,\mathrm{kg\,m^2}$. Rotational inertia describes resistance to changes in turning motion. That change belongs to the tested mechanism. No extra physical parameters were retuned and no alternative \texttt{implicitfast} fallback was used. This check does not calibrate the hand's position, orientation, size, contacts or fidelity to biology.

\subsubsection{Compare the two endpoints}

Use $P$ for the pulsed hand, $C$ for the control and $T=0.200\,\mathrm{s}$ for the nominal endpoint. The three reported differences are
\[
\Delta q_{\infty}=\max_{1\leq j\leq23}|q_{P,j}(T)-q_{C,j}(T)|,
\qquad
\Delta a=a_{P,\mathrm{FDS2}}(T)-a_{C,\mathrm{FDS2}}(T),
\]
\[
\Delta F=F_{P,\mathrm{FDS2}}(T)-F_{C,\mathrm{FDS2}}(T).
\]
The first calculation finds the difference between corresponding final joint angles, discards its sign and selects the largest of the 23 values. It describes the final instant, rather than the largest difference anywhere along the paths. The second subtracts the control activation from the pulsed activation for the selected muscle model. The third makes the same signed subtraction for its tendon-actuator force. Activation has no physical unit; force uses newtons and the model's sign convention. A forward evaluation recalculated each final force at its corresponding final state, avoiding force values left over from the last step.

% Shared by all four proposed R3 inserts; no added packages or font defaults.
\begin{table}[htbp]
\centering
\caption{Final contrasts for one simulated right-hand mechanism. All values compare the pulse branch with its matched control at nominal $T=0.200\,\mathrm{s}$. The angle value is the largest absolute separation across 23 hinges at that endpoint; activation and force retain the pulse-minus-control sign. Activation and force concern the selected \texttt{FDS2\_r} actuator. Force uses the synchronized final forward evaluations and the model's sign convention. These are simulated contrasts, not errors against biological measurements. The displayed digits preserve the recorded numerical values, without establishing biological accuracy.\cite{myohandNumericalRecord20261105}}
\label{tab:myohand-numerical-response}
\begin{tabular}{p{0.53\linewidth}p{0.26\linewidth}p{0.10\linewidth}}
\hline
Quantity & Contrast & Unit \\
\hline
Largest absolute final joint-coordinate separation, $\Delta q_{\infty}$ & $0.5027347840656727$ & rad \\
Selected final activation difference, $\Delta a$ & $0.05205729904534248$ & none \\
Selected final tendon-actuator force difference, $\Delta F$ & $-6.407722136774943$ & N \\
\hline
\end{tabular}
\end{table}


A separate standard-library audit checked the saved evidence without rerunning the model. It completed 132 checks without a failed check. The required numerical channels were finite, and recorded native warning counts stayed at zero. Four saved endpoint mass-matrix spectra contained positive values. A mass matrix describes the model's resistance to acceleration; its spectrum is a numerical diagnostic of that matrix. The raw matrices and the inputs to residual calculations were not independently exported. The audit therefore checked the stored diagnostics without independently solving those matrices again.

There was no observed contact. At the final control endpoint, the recorded soft-limit excess was $0.0004121829780273552\,\mathrm{rad}$. A soft limit resists crossing a boundary without proving an absolute hard stop. Neither energy conservation nor numerical convergence was established. The test imposed no physiological threshold for calling the response large enough.

The result is a difference between two simulated trajectories. It is not an error measured against a person's movement, muscle signals or forces. The 200 time steps depend on earlier states; they do not supply 200 subjects or 200 independent experiments. No statistical $p$-value follows. The decimal places preserve what the numerical record contains, without establishing biological accuracy or repeatability in other cases.

\subsubsection{Follow the evidence through to the trainee}

The recorded route used CPU dynamics. Its BLAS numerical-library pool reported one thread and its dynamics pool size was zero, while temporary parallel work during compilation was unmeasured. The 263 observed loaded-file identities cover that bounded inspection, not every dependency of the system. Resource use and speed remain unresolved. The process reported high-water memory of 5,187,464 KiB, a value that can retain pre-exec history; it cannot be assigned to the hand model alone or establish a model-only 3 GiB bound. Tool/native timing near 448 ms and timing inside the probe near 491.927 ms have unreconciled boundaries. They supply neither complete useful end-to-end latency nor a speed improvement, and they do not account for every CPU/GPU cost.

With zero sensors, the hand records none of Myoception's proposed sensory side of muscle-driven action, such as touch or body-position feedback. Calculated activation, force and motion are different from sensory nerve recordings or measured biological feedback. This test trains no controller or Large Movement Model and tests no generalization. The wider proposal still draws on four groups: neural observations, muscle measurements, measured movement and physics-generated examples. Before this generator could supply training examples, its mechanics, controls, calibration, uncertainty, source rights and data lineage would need qualification. Independent held-out real measurements would then test its contribution. More data and more varied data remain separate, untested hypotheses.

Even a later controller that passes its mechanical tests would leave Brain2Body's teaching question open. A trainee must be tested after trainer commands end, with any necessary compensatory movement bridge kept constant so that movement can still be expressed. Match the information, practice and assistance, retain the original sham, delayed and unrelated-guidance controls, and compare independent outcomes. No trainee, teaching, withdrawal or physical intervention appears in this diagnostic. It chooses no training generator and supplies no biological ground truth. LINK's eighteen fits and 720 saved epochs still fail both admission gates and all five benefit claims; protected tests and quarantined sources remain closed.

\subsection{Follow the gain until the teaching ends}

A better movement prediction might make assistance more useful. To test teaching, the experiment must also identify what the trainer supplied: expert information, feedback, commands, practice and the amount of help. Match those resources against an appropriate simpler controller and the original controls using sham, delayed or unrelated guidance. Otherwise, extra practice or being given the correct action could explain the apparent benefit.

The retention test stops trainer commands while keeping any necessary compensatory movement bridge constant. That bridge lets a trainee express movement despite an interrupted pathway; removing it together with teaching could make movement impossible regardless of what was learned. Immediate performance, later retention and unfamiliar variations are separate outcomes. A gain that requires continued guidance can support an assistance claim. A gain that persists under the controlled comparison needs its own uncertainty and tested domain. It still would not establish tissue repair or the ability to transfer any skill to any person.

Offline replay, simulation, measured biological feedback and physical intervention keep their own evidence and authority. The current proposal reports no scaling result, trained body-wide world model or biological teaching. LINK's eighteen fits and 720 saved epochs retain their failed model-admission gates and all five failed benefit claims. The next comparison has to show which prediction improves, what supplied the improvement and whether the trainee retains anything from it.

% Prospective assistance/retention protocol; no biological experiment reported.
\section{Catch it again when the help is gone}

A computer could guide a hand to catch a ball. The catch happened. What would
happen on the next attempt, after the teaching help stopped? That is the
Brain2Body question. Keep the task and the necessary access to movement
comparable, or the second attempt cannot tell us what the trainee retained.

The instruments introduced in \emph{From a reach to a testable model} would
have to agree on the catch. Scalp EEG records scalp electrical activity,
surface-EMG sleeves on both arms record muscle electrical activity, and the
synchronized, calibrated left/right tripod views would measure the hands and
arms. OpenBCI or an equivalent configurable system remains a candidate whose
sites, capacity, units, delays and clocks need qualification. The question now
is whether those observations can distinguish a successful catch from the
movement that merely looked close.

These left/right hand views differ from the earlier proposed top and bottom
dot maps. Both belong to the motion-capture relay-composite idea, but each
configuration needs its own calibration, timing and visibility checks. A point
the camera cannot see must not quietly become a measured point because a model
can draw it. The virtual body also needs appropriate anatomy and mechanics.
EMG is neither a force measurement nor a recording of incoming sensory nerves;
contact and force claims need observations that can check them.

\subsection{Let each action add a problem worth solving}

Begin by rubbing the hands rapidly together. Touch and movement happen at once.
Matched conditions would help tell their effects apart. The requested knuckle
movement and the stretch with interlaced fingers, palms outward and arms
extended would be recorded as distinct actions. The cold-hands example makes
rubbing familiar; it establishes no cold-exposure result or measured sensation.

Next, pass a juggling ball between the hands in the proposed hot-potato task,
then progress to the preferred two-ball juggling task. Release makes a new
problem for the catching hand. Check the model against the independently
observed handover and its timing, including what remains uncertain. A prediction
from the cue or previous movement might explain success without useful neural
information, so those comparisons stay in the test.

Boxing would ask the system to distinguish stated punch types from measured
movement. NeuroBoxer is the proposed virtual game built around that comparison.
Hands and wrists help stabilize punches as well as the arms moving. A prime-era
Mike Tyson opponent is a proposed virtual game character. The game supplies no
physical opponent or evidence about safe physical boxing. Expert footage can show
movement and tactics within its rights and camera limits. It cannot provide
that expert's brain activity, muscle activity or force measurements.

Toy archery would add another form of coordination: one hand holds while the
other draws, then the fingers release and the hands return. The proposal uses
arrows with small, cup-shaped suction tips and an over-arm pouch with several
arrows. NeuroArcher would use the coordination
against a virtual target. Hitting it cannot, by itself, establish what happened
in a biological nervous system.

The complete snap described in \emph{From a reach to a testable model} tests another
relationship: preparation, contact and return must all agree with the specified
movement. The cup reach remains the simpler precision comparison. Both need
organism-specific anatomy, contact, timing and real-reference evidence; a
human-looking animation cannot establish that a monkey can perform the snap.

\subsection{A virtual hand is the first control milestone}

The first game milestone would follow measured body motion. A later one would
ask whether permitted EEG observations can control a stated virtual action
without using withheld muscle or movement targets. Looking hands-free cannot
settle it. Eye activity, small covert muscle contractions, residual movement,
shared cues and future information need declared controls. Keep the original
brain-only, body-only and paired comparisons, including the simple prediction
that a target stays as it was, and keep the original success rules.

Understanding the intermediary should let us connect a stated intention with
an action as closely as possible for a defined task. The test needs a stated
person, condition, allowed context, tolerated error and independent observation.
It does not establish a universal one-to-one brain map. Predicting a recorded
EEG feature from body observations is likewise a conditional estimate. It does
not identify a unique mind or deliver a sensation.

\subsection{Teaching help and a movement bridge serve different roles}

Remove trainer information to test what was retained. Keep any movement bridge
needed to perform the task, and keep the comparison conditions meaningful.
Measure performance now, retention later and transfer to specified unfamiliar
conditions separately. Report how much help remains and what else could explain
the result. Removing vision in a proposed blindfolded condition changes feedback;
it does not make a test independent or prove learning by itself.

That retention comparison also depends on how movement or feedback reaches the
trainee. EEG and surface EMG supply recordings. A future body writer would need
a separate qualified interface: neuromuscular electrical stimulation (NMES)
and functional electrical stimulation (FES) can cause contractions or functional
movement, while a sensory writer or transducer has a different delivery role.
What the interface delivers would need its own measurement. Physical spinal
work remains read-only. These physical additions are still unqualified and
outside the acquisition and virtual-work budget scope; no such participant
recording, stimulation or actuation is reported here.

Training the computer model and teaching a biological trainee are different
questions. Qualified recordings and simpler baselines come before paired fitting
and fine-tuning. Reinforcement learning would need specified actions, resulting
transitions and rewards; a good reward inside a simulation is not a biological
teaching result. Final held-out observations stay outside training and selection,
and a frozen running model stays separate from a candidate adaptive update. The
current LINK comparison failed both model-admission gates and all five benefit
claims. This task sequence supplies future comparisons, not a new result about
retained skill, repaired tissue or a cure.

% Prospective apparatus and protocol; scientific and publication acceptance remain open.
\section{Build a brain--body system we can test}

What would justify building the next prototype? The paired observations in
\emph{From a reach to a testable model} and the tasks in \emph{Catch it again
when the help is gone} give that decision a test. The virtual hand would predict
from one declared input side; independent movement observations would check the
answer. A missed contact or a late movement identifies a problem to resolve
before expanding the apparatus.

The proposed design brings non-invasive EEG and surface EMG into one conditional
learner with the calibrated motion-capture relay-composite and an open
musculoskeletal-simulation physics-engine shown in a lifelike view. Its open
multimodal model would need eligible neural and muscular structure and activity
data. The EEG/EMG pair supplies intent-related brain and muscle observations;
the movement/physics pair tests what the body did and how the mechanics predict
it. Form constrains what could happen; measured function checks the prediction.
Neither pair can certify itself through agreement inside the learner.

These five roles give us brain recordings, muscle recordings, measured
movement, a computer body and a learner that relates them. The capture
composite would join checked movement views. The physics engine would
calculate movement under stated muscle, joint and contact rules. Multimodal
means that the learner can relate different kinds of observation; it does
not mean they were all recorded together or make a missing channel measured.

\subsection{One movement. Two directions.}

The two directional tests need different answers kept outside their inputs.
For the forward test, EEG and declared task context would predict supported
muscle activity and movement on withheld recordings. For the reverse test,
muscle activity and movement would predict recorded EEG features within the
task. Each prediction needs its own independent measured target; agreement
between models cannot substitute for either comparison.

Movement establishes what the body did, narrowing the plausible explanations
of the brain signals without identifying a unique intent or private experience.
Stimulus, task, behaviour and separately reported context therefore remain with
the observations. Uncertainty belongs beside the estimate; a generated neural
pattern cannot be displayed as the observed mind.

Surface EMG still records muscle electrical activity. An incoming sensory or
feedback-pathway claim needs its own evidence, which the virtual body cannot
supply. Recordings, predictions and simulation must remain identifiable as the
comparison follows one direction and then the other.

\subsection{A bridge must account for the spine}

The spinal cord helps shape the response. Muscle output depends on sensory
input, movement context, local circuitry and descending input; some spinal
pathways can retain learned changes. Reproducing a specified spinal function
would therefore require more than forwarding its signal: a brain--body bridge
must account for the processing it aims to reproduce.

A changed response can reflect immediate reflex modulation, lasting spinal
plasticity or peripheral adaptation. Those explanations need separate tests.
A change during the task does not establish retained memory, and a retained
change does not by itself locate that memory. Human reflex conditioning and
constrained mouse adaptation supply evidence for different claims.~\cite{thompson2009retained,thompson2013retained,lavaud2024}
The human studies do not establish brain-independent learning; the mouse
preparation does not demonstrate unrestricted walking or reaching. Our human
comparisons use abstracts and selected results, discussion or methods passages.
Relevant mouse-study passages were also examined in the author manuscript;
equivalence to the final published edition and the complete methods and
statistics have not been established here.

Reflex modulation means that the current conditions change a response;
plasticity means a change in the nervous system that can last. Peripheral
adaptation concerns changes outside the brain and spinal cord. A later
change in movement could involve more than one of these processes, which
is why persistence alone cannot tell us where or how the change occurred.

The periphery adds a different question: did the sensory signal itself change?
Passive proprioceptive training supplies a task-specific peripheral comparison;
its psychophysical and spindle-afferent experiments were separate.~\cite{ackerley2022retained}
Our passive-training comparison uses abstracts and selected results, discussion
or methods passages. None of these studies demonstrates universal skill
transfer or complex memory stored in a peripheral axon.

The psychophysical comparison tests what a person can distinguish or report.
The spindle-afferent comparison records sensory signals associated with
muscle spindles, receptors involved in sensing stretch. These were separate
experiments: placing their results together does not turn them into a single
simultaneous observation of perception and nerve activity.

The next comparison must choose a function we can measure. Does the controller
respond appropriately to the supported context? Does an adaptation persist
through the declared test? Could peripheral or mechanical changes explain the
outcome? Keep recordings, estimates and simulated state in distinct visible
layers. Surface EMG cannot supply a missing sensory-nerve or spinal channel;
an explanatory cutaway cannot register recorded neural sites.

Externalization remains the research aim: can a specified spinal function be
reproduced outside its usual pathway? Replacing a computation, stimulating a
pathway, controlling a simulated limb and restoring biological function would
answer different questions. Each needs its own task- and site-specific evidence;
physical spinal scope remains read-only. Anatomical access does not establish
neural function, increased excitability does not establish useful plasticity,
and no fixed percentage of a complex skill can be allocated to brain, spine
or periphery from these proposals.

\subsection{Reading a signal and writing an output}

A prediction leaves the delivery question open: what, if anything, reached
the trainee? \emph{Write sensation. Enact action} distinguishes the computer's
input/output boundary from afferent/efferent direction relative to the central
nervous system; a robot command has an action role without becoming an
anatomical efferent nerve. That distinction determines the further test.

Delivering sensation would require a separately qualified sensory writer or
transducer. Indirect sensory feedback and direct neural writing need distinct
device, measurement, safety and clinical evidence. Producing an action needs
the appropriate virtual, robotic or biological output interface. The recording
prototype does neither by predicting EEG or EMG. Physical sensory writers,
transducers and body actuators remain excluded from the acquisition and
virtual-work plan, unqualified and unpriced. This proposal authorizes no
stimulation, participant acquisition, biological actuation, new equipment
purchase or training.

\subsection{The first prototype has to earn the next one}

The lean tier would first ask whether the selected signals and calibrated
movement reference can be synchronized for a bounded task. The serious-prototype
tier would improve acquisition, independent references and source-specific
bidirectional evaluation. Evidence from those comparisons would determine
whether the research-lab tier should expand cohorts, tasks, instrumentation and
compute. A larger apparatus needs an answered question from the smaller one.

Each tier needs equipment, calibration, accessories, storage, integration,
appropriate safety and consent work, independent references and contingency.
Current vendor sources should distinguish listed prices from engineering
allowances and items needing quotes. Planning figures stay in USD, with a
dated basis for any currency conversion. These budgets describe proposed
scopes; they authorize no purchase and promise no complete neural reconstruction.

\subsection{Let the model earn the claim}

The first model comparison depends on compatible observations, not the
length of the source catalogue. Rights, version, coverage, clock, units and
missing modalities determine what an open dataset can contribute. Fitting
would require an independently admitted protocol, with participants, sessions
and tasks excluded from model selection reserved for testing. Simple baselines
would show whether the added model helps; artifact leakage, absent inputs,
timing errors, physically invalid solutions and uncertainty could explain where
it fails. Earlier failures and reserved recordings remain part of that test's
constraints.

Constrained virtual hand or foot control would be the starting domain. Robotics,
remote operation, unfamiliar effectors and extra limbs remain later research
directions, each requiring separate physical, latency and safety evidence.
A successful movement prediction supports its movement claim. Biological and
perceptual claims need the corresponding broader evidence.

Rights-cleared images and lifelike interactive models would let a reader inspect
the apparatus and compare the observation, prediction and error under their
declared conditions. A convincing image cannot fill a missing measurement or
demonstrate biological function. The first prototype would have to make the
difference visible and testable before the next one could rely on it.


\section{Measured Improvement}

The programme can begin with the key press: what should move, and what result would count as useful? Formulating means deciding that question; specifying means naming observable outcomes that would answer it. Verification then checks whether the implementation does what was intended. Evaluation compares its answers with independent references. Measured error shows what to refine next.\par

These parts need not happen only once. A failure may expose a poor question, a missing measurement or an incorrect implementation. The response should diagnose the end-to-end cause while keeping the original evidence and thresholds. Repeating work for the sake of activity does not create scientific progress.\par

Model size and computing speed also need evidence. A larger model may help some tasks; it cannot create missing paired measurements. Faster training on a machine is a computer-workload result, not a promise of faster biological mastery. A hardware change earns promotion only through matched useful outcomes with the same quality and acceptance requirements.\par

Generated data can help practice and controlled development tests. It cannot supply independent measured biology for a target never observed. Real held-out measurements, comparison methods, uncertainty and failures remain necessary. Savings in time or effort are useful only when the scientific question and all its gates remain intact.\par

\section{Internal-Body Sensing}

Interoceptomics is the programme's proposed study of internal-body sensing together with maps of the neural connections involved. Heartbeat, breathing and other internal signals could offer additional questions beyond visible movement. The word combines interoception, sensing the body's internal condition, with connectomics, studying neural connections.\par

Those questions need appropriate observations and outcomes. A finger-contact sensor does not directly measure the feeling of a heartbeat. A model of breathing motion does not by itself reveal private sensation. The movement-first strategy can help ground some relationships without turning every experience into an objective readout.\par

The names NeuroTranslation and Interoceptomics describe proposed scope. NeuroTranslator describes implemented software. NeuroTranslation Network describes a proposed design. Brain2Body names the overall programme. NeuroTranslate names a proposed mapping operation with declared inputs, targets and a task. Testing one direction does not establish its reverse, a unique brain state or delivered sensation. Keeping these meanings consistent prevents a new label from being mistaken for a new demonstrated capability.\par

\section{Tests That Could Reject the Proposal}

The framework loses support where appropriate controls do just as well, the mapping fails on later recordings, or trainer gains disappear after help stops. A simulation that violates independently supported physical bounds cannot support the claimed biological mechanism. A claim whose target cannot yet be measured remains unvalidated rather than being promoted through visual confidence.\par

Negative results are not a reason to conceal the ambition. They tell us which next question is justified. The current development models fail their original benefit gates. The older classifier, keyboard replay and goal-driven policy remain separate bounded results. None demonstrates arbitrary skill teaching, complete nervous-system coverage or a voluntary central--spinal--peripheral chain.\par

There may still be useful limited interfaces even if a universal translator never appears. A narrowly tested prediction or a carefully measured aid can matter. Its value comes from what it demonstrably helps someone do under stated conditions, not from claiming every possible future result at once.\par

Return to the tune. If the hand needs teaching commands for every note, the
result may still be useful assistance. If the trainee performs better later
without those commands, while retaining any necessary movement bridge, we
have a result worth examining for retained learning. We would still need the
matched controls, uncertainty and unfamiliar variations to say where that
gain holds. Brain2Body's next generation must show what moves, who or what
supplied the help, what remains when teaching ends and where the explanation
breaks. This article proposes that comparison. It does not claim the trainee
has already passed it.\par

\section{Definitions and Lexicon}

This \hyperref[programme-definitions]{Simple English Edition definition and lexicon supplement} uses the same meanings as the scientific article. It is derived from the shared Brain2Body registry identified by version and fingerprint at the start of that supplement. A fingerprint identifies a file; it is not a scientific score. Alternative historical labels remain lookup aids. A glossary supports the explanations in the text rather than replacing them.\par

\section{References and Evidence Status}

The original scientific sources are linked beside the statements they support.
Dated developer accounts and pinned code observations are identified separately
from published biological experiments. The accompanying
\href{research-evidence.json}{computational results and limits} describe the
programme's unpublished numerical and software studies. The companion
NeuroTranslation paper is an unpublished working draft. The historical
conceptual figure explains a relationship; it supplies no new biological
evidence. The edition note records the remaining scientific and publication
limits.\par

